\documentclass[11pt]{amsart}
\usepackage[letterpaper,margin=1in]{geometry}
\usepackage[T1]{fontenc}
\usepackage{lmodern,microtype,amsmath,amssymb,amsthm,mathtools,mathrsfs}
\usepackage{booktabs,array,enumitem}
\usepackage{graphicx}
\usepackage[section]{placeins}
\usepackage{tikz}
\usetikzlibrary{arrows.meta,positioning,calc}
\usepackage[hidelinks]{hyperref}
\hypersetup{pdftitle={An embedded band foliated by closed asymptotic lines and nonrigid smooth tight tori}}
\newtheorem{theorem}{Theorem}[section]
\newtheorem{proposition}[theorem]{Proposition}
\newtheorem{lemma}[theorem]{Lemma}
\newtheorem{corollary}[theorem]{Corollary}
\theoremstyle{definition}
\newtheorem{definition}[theorem]{Definition}
\theoremstyle{remark}
\newtheorem{remark}[theorem]{Remark}
\newcommand{\R}{\mathbb R}\newcommand{\Sph}{\mathbb S}
\newcommand{\dd}{\,\mathrm d}\newcommand{\Id}{\mathrm{Id}}
\newcommand{\e}{\mathrm e}\newcommand{\ee}{e_3}
\newcommand{\ip}[2]{\langle #1,#2\rangle}\newcommand{\norm}[1]{\lVert #1\rVert}
\DeclareMathOperator{\supp}{supp}
\DeclareMathOperator{\conv}{conv}\DeclareMathOperator{\Stab}{Stab}


\DeclareMathOperator{\grad}{grad}
\DeclareMathOperator{\Lk}{Lk}
\DeclareMathOperator{\Isom}{Isom}\DeclareMathOperator{\Emb}{Emb}
\newcommand{\II}{\mathrm{II}}\newcommand{\Gmap}{\mathcal G}


\numberwithin{equation}{section}
\setlist{nosep,leftmargin=1.7em}
\setcounter{tocdepth}{2}
\emergencystretch=1.5em
\title[Closed asymptotic bands and tight tori]{An embedded band foliated by closed asymptotic lines\\and nonrigid smooth tight tori}
\author{}
\date{}
\subjclass[2020]{Primary 53A05; Secondary 53C24}
\keywords{Closed asymptotic lines, infinitesimal bending, tight torus, differential Legendre transform, radial saddle completion, complete saddle cylinder, asymptotic holonomy, rigidity}

\begin{document}
\begin{abstract}
We construct a smooth embedded negatively curved band foliated by closed
asymptotic curves. The full return map of its nonruling asymptotic family
is the identity, and its compactly supported infinitesimal bendings are
parametrized by arbitrary smooth functions of a transverse leaf coordinate.
We complete the band to a smooth embedded tight torus while retaining
its bending region. Opposite signs of one retained bending give two
noncongruent tight embeddings with exactly the same intrinsic metric;
they agree on a nonempty open set and have no open planar subset.
Countably many disjoint bendings give a Cantor family with these properties.

The construction separates spatial closure from cancellation of the
asymptotic return period. A corrugated spherical normal loop and a
positive speed provide an embedded identity-return band with two visible
exits. Visible connectors, differential Legendre duality, and radial
saddle filling then produce explicit end profiles. Compactification and
convex attachment close the surface, and an affine marking removes its
ambient symmetries. The same band also gives fixed-boundary isometric
saddle annuli, an annulus with injective Gauss map and zero normal
linking along its closed asymptotic curves, and a complete properly
embedded negatively curved cylinder with total absolute curvature $4\pi$.
These constructions connect the smooth tight-congruence question with
the geometry and return dynamics of closed asymptotic curves.
\end{abstract}
\maketitle
\section*{Formalization objective and route}
The first objective is the two-realization assertion of Theorem~\ref{thm:main-fiber}. Cantor-family, complete-cylinder and linking applications remain deferred. The first-pair construction status below is derived from exact current kernel declarations; the canonical final theorem has no background or construction premises.
The CURRENT exact combined kernel report certifies every original construction gate and the ultimate first-pair theorem. The first-pair objective is unconditionally proved. Positive-Gauss tightness, ellipse recognition, equal-image reparameterization and fixed-open rigidity have exact closed kernel proof inhabitants. The canonical theorem has no background or construction premises. SAME seed, FINAL Y, clocks, raw source E, one full-domain H with its OWN gradient inverse, completed scalar/beta/Q/Cdata tuple and ONE FULL meridian are retained. Full-paper and extension/Cantor-family scope remain incomplete or deferred.


\section{Introduction}\label{sec:intro}
\subsection{Congruence, tightness, and the historical problem}
When does the intrinsic metric of a surface determine its image in
$\R^3$ up to a Euclidean motion? Convex-surface rigidity provides the
classical model, but a surface with both positive and negative curvature
poses a different problem. The negative-curvature region supports a
hyperbolic system, and its asymptotic curves are the characteristics.
Yau emphasizes this connection between metric uniqueness and asymptotic
curves in his review of the global embedding problem
\cite[pp.~235--236]{YauReview}.

Tight surfaces form a natural class in which to test uniqueness beyond
convexity. For a closed embedded surface, tightness is equivalent to the
two-piece property and to attaining the topology-dependent lower bound
for total absolute curvature. A tight torus has total absolute curvature
$8\pi$. Aleksandrov proved congruence for analytic tight tori
in 1938. Nirenberg's 1963 work obtained a smooth partial result.
Its torus formulation assumes $C^5$ regularity and $\nabla K\ne0$ at
$K=0$, together with an additional asymptotic-curve hypothesis. The two
torus versions are recorded together in \cite[Theorem~1.6.1]{BK}; the original
smooth argument is \cite{Nirenberg}.

The distinction between smooth and polyhedral surfaces is already
present in the earlier literature. Banchoff constructed noncongruent
isometric tight polyhedral tori; Banchoff--K\"uhnel record these examples
and then ask whether two isometric smooth tight tori must be congruent
\cite[Example~1.6.2 and Question~11]{BK}. Ghomi's later survey states
the smooth tight-rigidity problem and its fixed-boundary
negative-curvature annulus counterpart \cite[Problems~1.2--1.3]{GhomiProblems}.
Thus the smooth tight-congruence question has a specific history;
nonuniqueness for arbitrary smooth surfaces is an older phenomenon.

Han--Khuri subsequently proved two rigidity theorems that allow several
closed asymptotic curves. Besides their hypotheses at the boundary of
the negative-curvature region, both theorems require
\[
 \oint_\Gamma k_gk_n|K|^{-1/2}\,ds\ne0
\]
on every closed asymptotic curve; $k_n$ is the normal curvature in the
tangent direction perpendicular to $\Gamma$. They also ask whether
tightness itself forces this nonvanishing \cite[Theorems~1--2 and the
following remark]{HK}. Our protected band has identity return, so this
integral vanishes on every one of its leaves. The resulting example
therefore lies outside the nondegeneracy hypothesis of those theorems.

Our main result gives two smooth embedded tight tori with one common
metric, noncongruent images, and an open region of pointwise agreement.
The entire change is supported inside a compact negative-curvature
region. We also obtain a Cantor family of such realizations.
Here two embedded surfaces are \emph{congruent} if an ambient Euclidean
isometry carries one image onto the other. The following statement therefore
rules out congruence of images, as well as congruence of parametrized maps.
Write $g_{\rm Euc}$ for the Euclidean metric.

\par\noindent\textbf{Formalization: pending.} See the precise object and dependency interfaces in the blueprint.\par
\begin{theorem}[Localized nonuniqueness for tight tori]\label{thm:main-fiber}
There are a smooth metric $g$ on $T^2$ and two smooth tight embeddings
$X_+,X_-:T^2\to\R^3$ such that
\[
 X_+^*g_{\rm Euc}=X_-^*g_{\rm Euc}=g.
\]
Their images are noncongruent, they agree on a nonempty open subset, and
neither has an open planar subset.

In fact, the same metric can be chosen to admit a continuous family
\[
 \varepsilon\longmapsto X_\varepsilon\in\Emb_g(T^2,\R^3),
 \qquad \varepsilon\in\{-1,1\}^{\mathbb N},
\]
of tight embeddings with pairwise noncongruent images. All members agree
on the same nonempty open set, have no open planar subset, and have total
absolute curvature $8\pi$. The induced map into
$\Emb_g(T^2,\R^3)/\Isom(\R^3)$ is a topological embedding.
\end{theorem}

Here $\Emb_g(M,\R^3)$ is the space of smooth embeddings inducing $g$, with
the $C^\infty$ topology; $\Isom(\R^3)$ acts by postcomposition. The sign
space has its product topology and is homeomorphic to the Cantor set.
The open set in the theorem is fixed by the construction. The family
is continuous on a Cantor parameter space, which is totally disconnected.
It supplies no nonconstant path of isometric embeddings. In particular,
the theorem does not settle the continuous smooth-flex problem stated
in \cite[Problem~1.1]{GhomiProblems}.

\subsection{Why localized bendings give exact pairs}
The exact pair comes from a quadratic identity. For an immersion $X$ and a vector
field $y:M\to\R^3$, the metric satisfies the exact identity
\begin{equation}\label{eq:intro-quadratic}
 g_{X+y}=g_X+S_Xy+\langle dy,dy\rangle,
\end{equation}
where $S_Xy$ is the linearized metric change. If $S_Xy=0$, then
$g_{X+y}=g_{X-y}$. Thus a nonzero infinitesimal bending gives an exact
sign ambiguity, provided both perturbed maps remain embeddings. Their
common metric is generally different from $g_X$.

If finitely many such fields have disjoint supports, then
\[
 g_{X+\sum_{j=1}^m\varepsilon_j a_jy_j}
 =g_X+\sum_{j=1}^m a_j^2\langle dy_j,dy_j\rangle,
 \qquad \varepsilon_j\in\{-1,1\}.
\]
This is the finite sign mechanism. The Cantor family follows by choosing
countably many disjoint fields and rapidly decreasing amplitudes. The
central issue is therefore geometric: construct a tight torus carrying
localized bendings in a compact negatively curved region, and eliminate
ambient symmetries on the region that stays fixed.

\subsection{The band that carries the construction}
Our bending region lies in a principal-normal ruled band
\[
 X(s,u)=\gamma(s)+uE(s).
\]
Here $s$ is arclength along a closed curve, and the periodic frame
$(T,E,n)$ satisfies $T_s=kE$, $E_s=-kT+\tau n$, and $n_s=-\tau E$,
with $\tau\ne0$. The band has strictly negative Gaussian curvature.
Put $\rho=\sqrt{|\tau|}$ and $D=k_s-k\tau_s/\tau$. Its nonruling
asymptotic curves have return map
\begin{equation}\label{eq:intro-return}
 \mathcal R(u)=\frac{u}{1+\rho(0)I_0u},
 \qquad I_0=\oint\frac{D}{2\rho}\,ds.
\end{equation}
The multiplier $\mathcal R'(0)$ is always one. The stronger condition
$I_0=0$ makes the full return map the identity, so nearby characteristics
are closed and admit a single-valued transverse parameter.
Theorem~\ref{thm:ruled} proves that this condition is exactly what permits
compactly supported infinitesimal bendings on a thin one-sided band. In
the identity case, every such bending comes from one arbitrary transverse
test function.

Producing that band and completing it are separate tasks. For a prescribed
spherical normal loop, spatial closure is a vector moment condition,
whereas identity return is one scalar period condition.
Theorem~\ref{thm:normal-loop-criterion} separates them and cancels the
period without losing closure. The explicit corrugated seed supplies
additional embedding and two-sided visibility properties. We then protect
a compact subannulus, alter the support only outside it, and attach the
remaining geometry.


Figure~\ref{fig:band} shows the thin balanced band in the saved numerical
atlas and a local view of its protected region. On the unchanged region, its rescaled leaf label
$c$ is related to the first integral used below by
$v=(\sigma c)^{-1}$, with $\sigma=10^{-6}$ in this numerical model.
The unrolled strip resolves leaves that cannot be separated in the full
physical view. The explicit analytic seeds used in the proof are drawn
separately in Figure~\ref{fig:seeds}.
\begin{figure}[htbp]
\centering
\includegraphics[width=\textwidth]{figures/01-band.pdf}
\caption{The balanced numerical band used by the atlas's saved completion.
(a) The complete band, with a protected closed leaf in blue. Its actual
width is too small to resolve at this scale.
(b) A rigidly rotated, uniformly magnified local patch: one displayed
length unit equals $10^{-6}$ original units. Gold marks the unchanged
closed-leaf region; blue and teal curves are its asymptotic leaves, and
the short transverse lines are straight rulings there.
(c) An unrolled parameter diagram, rather than a physical projection,
with protected levels $0.35\le c\le0.65$ and the two ends identified.
This saved balanced numerical example illustrates the construction;
the analytic existence argument is Lemma~\ref{lem:seed} and
Theorem~\ref{thm:spike}.}
\label{fig:band}
\end{figure}


\subsection{Further consequences and earlier questions}
The same protected band has two useful completions besides the tight torus.
\par\noindent\textbf{Formalization: pending.} Deferred from the first-pair objective.\par
\begin{theorem}[Closed asymptotic annulus]\label{thm:main-annulus}
There is a smooth embedded compact annulus $N\subset\R^3$ with $K<0$
in its interior and circular boundary components lying in tangent planes.
For a suitable normal, its Gauss map is a diffeomorphism
\[
 N^\circ\longrightarrow\Sph^2\setminus\{\pm e_3\}.
\]
An open subannulus is foliated by embedded closed asymptotic curves.
Each has linking number zero with a sufficiently small normal push-off.
\end{theorem}

\par\noindent\textbf{Formalization: pending.} Deferred from the first-pair objective.\par
\begin{theorem}[Complete embedded saddle cylinder]\label{thm:complete-cylinder}
There is a smooth proper embedding $X_\infty:S^1\times\R\to\R^3$
whose image $\Sigma$ has complete induced metric and $K<0$ everywhere.
It contains an open annulus of embedded closed asymptotic curves with
identity return and zero normal linking. A choice of normal gives
\[
 n:\Sigma\longrightarrow\Sph^2\setminus\{\pm e_3\}
 \quad\text{a diffeomorphism},\qquad
 \int_\Sigma|K|\,dA=4\pi.
\]
For some nonempty interval $I$, there is an injective linear family of
infinitesimal bendings parametrized by $C_c^\infty(I)$, all supported in
one fixed compact subset of $\Sigma$.
\end{theorem}

The relation to the published questions is as follows.
\begin{enumerate}[label=\textup{(\roman*)}]
\item Da Cruz--Garcia record the question of a cylindrical region
foliated by closed asymptotic lines \cite[p.~17]{DaCruzGarcia}.
The identity-return band provides such a region in a surface in $\R^3$.
Their paper studies asymptotic lines of general plane fields; our
construction is on an actual embedded surface.
\item Nirenberg's annulus question concerns closed asymptotic curves
on a negatively curved annulus with convex planar boundaries, as
explained in \cite[\S1.2]{GhomiProblems}. Theorem~\ref{thm:main-annulus}
provides an open band of closed leaves in an annulus of this type.
\item Ghomi--Raffaelli ask whether an embedded negatively curved
surface with injective Gauss map can have a closed asymptotic curve
with zero normal linking \cite[Problem~1.1]{GR}. Our annulus has both
properties. The normal along a leaf is a continuous generalized
binormal, so it also gives the example asked about in
\cite[Problem~1.4]{GhomiProblems}. Here inflections are permitted;
a classical Frenet binormal need not be defined at every point.
Ding's recent open-annulus construction also proves the zero-linking
and regular-binormal conclusions \cite[Theorem~1.1 and
Corollary~1.2]{Ding2026}. The boundary geometry and closed-leaf
foliation in Theorem~\ref{thm:main-annulus} give additional conclusions.
\item Theorem~\ref{thm:main-fiber} gives a negative answer to the
smooth tight-torus congruence question \cite[Question~11]{BK}.
Corollary~\ref{cor:fixed-boundary-annuli} gives noncongruent isometric
negative-curvature annuli that agree on full boundary collars,
addressing \cite[Problem~1.3]{GhomiProblems}.
\item The continuum of closed leaves shows that the closed-curve
restriction discussed by Yau \cite[p.~236]{YauReview} does not hold
for all smooth tight tori. Corollary~\ref{cor:HK-degenerate} gives the
vanishing integral relevant to the remark following Han--Khuri's
rigidity theorems \cite[p.~464]{HK}.
\item Ni--Zhang--Zhang construct a complete negatively curved immersed
cylinder containing a closed asymptotic curve and explicitly leave
embedded existence open \cite[Theorem~3.4 and the following paragraph]{NiZhangZhang}.
Theorem~\ref{thm:complete-cylinder} provides a proper embedded example.
\end{enumerate}

These are questions and formulations in the cited works. The geometry
needed to meet their hypotheses is proved below: embedding and global
Gauss-map injectivity, full boundary collars, zero linking, and
completeness are separate conclusions of the construction.

\subsection{Relation to earlier constructions}
Closed asymptotic curves and principal-normal ruled strips are established
parts of surface theory. Kovaleva's 1968 construction gives a negatively
curved tube-type surface containing a closed asymptotic line
\cite{Kovaleva}; its role in the principal-normal strip construction
is recalled in \cite[Note~2.1]{GR}. Garcia--Sotomayor study the qualitative behavior and
structural stability of periodic asymptotic lines
\cite{GarciaSotomayor}. Our use of principal-normal rulings builds on
this earlier setting. The additional local requirement is identity of
the full return germ, which gives a whole band of closed leaves and
allows us to classify its compactly supported bending fields.

The topological restrictions are equally important. Ghomi--Raffaelli
relate normal linking to the writhe and rotation of projections of
closed asymptotic curves. They discuss earlier examples with zero
linking and noninjective Gauss map, and construct examples with
injective Gauss map and nonzero linking \cite[\S\S1,4]{GR}.
Their Problem~1.1 asks whether zero linking and Gauss-map injectivity
can occur together. Their analysis also explains why injective normals
force inflections \cite[Note~5.1]{GR}. This is why the signed frame in
our ruled band allows $k=0$: requiring a classical Frenet frame
everywhere would exclude the geometry under consideration.

Smooth isometric nonuniqueness also predates the present construction.
Sa Earp--Toubiana give closed smooth surfaces of arbitrary genus with
nondiscrete isometric realizations, including families parametrized by
nowhere dense compact sets \cite[Theorem~2]{SaEarpToubiana}. Their
construction reflects pieces across planar seams at which the profile
is flat to every order. A Cantor set of realizations is therefore not
by itself a new phenomenon. Nor does a planar seam imply an open
planar patch; the distinction between these conditions matters when
comparing examples.

Sabitov constructs two noncongruent isometric negatively curved
surfaces with a common open set \cite{Sabitov2018}. His later work
describes metrics of revolution with a discontinuum of noncongruent
isometric immersions \cite{Sabitov2021}. Local agreement and large
sets of isometric realizations thus have their own earlier examples.
The contribution here is their realization inside smooth embedded
\emph{tight tori}, with the changes supported in strict negative
curvature, together with a band whose full asymptotic return and
supported bending space are explicitly controlled. The absence of
open planar patches is one of the proved geometric properties;
it is not used as a blanket distinction from all earlier constructions.

Ding's October 2026 preprint constructs an embedded open negatively
curved annulus with injective Gauss map and a closed asymptotic curve
of zero normal self-linking \cite[Theorem~1.1]{Ding2026}. It also gives
a regular injective continuous binormal \cite[Corollary~1.2]{Ding2026}.
Its planar curves are obtained by simultaneous polygonal rounding and
mixed-integral cancellation; the theorem makes no completeness claim
and prescribes no tangent-plane boundary geometry. The analytic seed
used here supports a different continuation: we cancel the full return
period, classify the supported kernel, and retain a closed-leaf band in
the compact annulus, tight torus and complete cylinder. The overlapping
open-annulus and binormal conclusions are therefore also stated in
this recent preprint, alongside the older results described above.

\subsection{Proof architecture}
The construction has two layers. First we manufacture the localized
bending region. Then we complete and use it. The dependencies are:
\[
 \begin{gathered}
 \text{normal loop and closing speed}
 \ \longrightarrow\ \text{period balance}
 \ \longrightarrow\ \text{identity-return ruled band},\\
 \text{corrugated seed}
 \ \longrightarrow\ \text{embedding and two-sided visibility},\\
 \text{protected band}
 \ \longrightarrow\ \text{positive exits}
 \ \longrightarrow\ \text{circular gradient seams},\\
 \text{Legendre duality and radial filling}
 \longrightarrow\ \text{parabolic end collars},\\
 \text{saddle annulus}
 \ \longrightarrow\ \text{asymmetric tight closure}
 \ \longrightarrow\ \text{exact sign pairs and families}.
 \end{gathered}
\]
Section~\ref{sec:theory} proves the local kernel theorem;
Section~\ref{sec:structural-balance} handles closure and period balance.
Section~\ref{sec:protected-core} uses the seed to build the embedded core,
and Section~\ref{sec:exits} creates its positive exits.
Sections~\ref{sec:completion} and~\ref{sec:model-saddle} complete the
support and assemble the tight torus. Section~\ref{sec:reconstruction}
first proves finite sign branching, then its countable extension.
Section~\ref{sec:applications} proves the stated consequences.
The general projective ruled theory and the quantitative seed estimates
are in Appendices~\ref{sec:general-ruled} and~\ref{app:seed-estimates};
neither interrupts the main construction.

\section{Conventions and geometric tools}\label{sec:prelim}
\subsection{Metrics, bendings, and tightness}
All maps are smooth unless stated otherwise. A compact annulus includes
its boundary; $A^\circ$ denotes its interior. For an open interval or
surface, $C_c^\infty$ denotes smooth functions with compact support.
Writing $V\Subset U$ means that the closure of $V$ is compact and
contained in $U$. The torus $T^2$ is $S^1\times S^1$, with
$S^1=\R/2\pi\mathbb Z$.
Norms on compact sources use fixed smooth atlases or a fixed background
metric. For an oriented immersion $X$ with unit normal $N$, we use
\[
 g_X(v,w)=\langle dX(v),dX(w)\rangle,\qquad
 \II_X(v,w)=\langle D_v(dX(w)),N\rangle.
\]
Here $D$ is the Euclidean connection, applied to the ambient vector
field $dX(w)$; the displayed normal component is independent of the
extensions of $v,w$. The shape operator is defined by
$g_X(A_Xv,w)=\II_X(v,w)$, so $dN(v)=-dX(A_Xv)$ and $K_X=\det A_X$.
Changing the normal changes the sign of $\II_X$ and $A_X$, and leaves
$K_X$ and the asymptotic directions unchanged. The strain is the symmetric tensor
\[
 S_Xy(v,w)=\langle dX(v),dy(w)\rangle+
             \langle dX(w),dy(v)\rangle.
\]
An \emph{infinitesimal bending} is a field $y:M\to\R^3$ with $S_Xy=0$;
this definition allows both tangential and normal components. We write
$\Gmap(X)=g_X$ for the metric reconstruction map. The tensor
$\langle dy,dy\rangle$ means $(v,w)\mapsto\langle dy(v),dy(w)\rangle$.

An asymptotic direction is a nonzero vector $v$ with $\II_X(v,v)=0$.
When $K<0$, these directions form two transverse smooth line fields.
A \emph{return map} follows one such field once around a closed leaf,
starting and ending on a short transverse arc. Identity return means
identity of the whole germ, rather than only its first derivative.

A closed connected embedded surface is \emph{tight} if its intersection
with each open half-space is connected whenever nonempty. Its total
absolute curvature satisfies \cite[\S1.1]{BK}
\begin{equation}\label{eq:tight-bound}
 \int_M|K|\,dA\ge2\pi(4-\chi(M)),
 \qquad\text{with equality exactly for tight embeddings}.
\end{equation}
For a torus the equality value is $8\pi$.

We repeatedly use openness of compact embeddings in the $C^1$ topology.
Local projections to tangent planes exclude collisions between nearby
source points, uniformly on a finite cover; compactness separates the
remaining pairs. The same argument shows that an immersion injective on
a compact embedded curve is injective on a sufficiently thin neighborhood
of that curve.

For an embedded closed curve $\Gamma$ and a nonzero field $\nu$ transverse
to its tangent, define
\[
 \Lk(\Gamma,\nu)=\Lk(\Gamma,\Gamma+\epsilon\nu)
\]
for small $\epsilon>0$. We compute linking by intersection with an
oriented spanning surface \cite{Rolfsen}.

\subsection{Support coordinates}\label{sec:support}
The completion is easiest to control by using the normal as the source
coordinate. Let $q\in\Omega\subset\Sph^2$, with round metric $g_0$, and
let $H:\Omega\to\R$. Define
\begin{equation}\label{eq:support}
 a_H=\nabla H+Hq,\qquad B_H=\nabla^2H+H\Id,
 \qquad R_H=\det B_H.
\end{equation}
The gradient and Hessian in this formula use the round metric $g_0$.
In the sum with $H\Id$, $\nabla^2H$ denotes the endomorphism obtained
by raising one index of the covariant Hessian. Equivalently, the
associated bilinear form is $\operatorname{Hess}_{g_0}H+Hg_0$.
We write $B_H(v,w)=g_0(B_Hv,w)$ when a bilinear evaluation is needed.
The determinant $R_H$ is computed in a $g_0$-orthonormal frame.

\par\noindent\textbf{Formalization: proved.} See the precise object and dependency interfaces in the blueprint.\par
\begin{proposition}[Support reconstruction]\label{prop:support}
The $\R^3$-valued differential of $a_H$ is $B_H$. Where $B_H$ is
invertible, $a_H$ is an immersion with normal $q$, metric
$g_{a_H}(v,w)=g_0(B_Hv,B_Hw)$, second fundamental form
$\II_{a_H}(v,w)=-g_0(B_Hv,w)$, and curvature $R_H^{-1}$.
The shorthand $B_H^2$ for its metric refers to this bilinear form. Every immersion
with an injective regular Gauss map has this representation on its Gauss
image.
\end{proposition}
\begin{proof}
Differentiating $\nabla H$ in $\R^3$ produces the normal term
$-dH\,q$, which cancels the corresponding term from $d(Hq)$.
The remaining differential is $\nabla^2H+H\Id$. It is tangent to
$q^\perp$, so $q$ is normal. Differentiating $da_H(w)\cdot q=0$ gives
$\II(v,w)=-\langle B_Hw,v\rangle$. Thus the shape operator is
$-B_H^{-1}$ and its determinant is $R_H^{-1}$.
Conversely, in Gauss coordinates set $H=X\cdot q$. Then
$dH(v)=X\cdot v$, because $dX(v)\cdot q=0$, and the tangential and normal
parts of $X$ are $\nabla H$ and $Hq$.
\end{proof}

A self-adjoint field is \emph{Codazzi} when
$(\nabla_vB)w=(\nabla_wB)v$ for all tangent vectors $v,w$.
\par\noindent\textbf{Formalization: pending.} Deferred from the first-pair objective.\par
\begin{proposition}[Codazzi realization and periods]\label{prop:codazzi}
For a smooth self-adjoint field $B$ on $\Omega$, the vector-valued form
$\beta_B(v)=Bv$ is closed exactly when $B$ is Codazzi. On a simply
connected domain this is equivalent to $B=B_H$, with $H$ unique modulo
ambient linear functions. On a general domain one must also require
\[
 \oint_\gamma\beta_B=0\quad\text{for every closed curve }\gamma.
\]
\end{proposition}
\begin{proof}
The normal part of $d\beta_B(v,w)$ is
$-\langle v,Bw\rangle+\langle w,Bv\rangle=0$ by self-adjointness.
Its tangential part is $(\nabla_vB)w-(\nabla_wB)v$, the Codazzi defect.
If the form has a primitive $a$, put $H=a\cdot q$. Since $da$ is tangent,
$dH(v)=a\cdot v$, hence $a=\nabla H+Hq$ and $B=B_H$.
The primitive is unique up to a constant vector, which changes $H$ by a
linear function. Vanishing periods are precisely the obstruction to a
global primitive.
\end{proof}

On the northern hemisphere introduce \emph{gnomonic coordinates}
\[
 q(z)=\frac{(z,1)}{w(z)},\qquad w(z)=\sqrt{1+|z|^2},
 \qquad G(z)=w(z)H(q(z)).
\]
In these planar source coordinates the immersion is
\begin{equation}\label{eq:planar}
 X_G(z)=\bigl(\nabla G(z),G(z)-z\cdot\nabla G(z)\bigr),
 \qquad K_{X_G}=\frac1{w^4\det D^2G}.
\end{equation}
Indeed, writing $A=D^2G$ gives
\[
 dX_G(\xi)=(A\xi,-z\cdot A\xi),\qquad
 g_{X_G}=A^T(I+zz^T)A,\qquad \II=-A/w.
\]
Thus $R_H(q(z))=w^4\det D^2G$. A \emph{saddle support} means
$R_H<0$, equivalently $\det D^2G<0$. Planar gradients below are
Euclidean, in contrast with the spherical gradient in \eqref{eq:support}.

If $p$ is a planar source curve, let $\gamma=\nabla G\circ p$ be its
gradient trace and $g=G\circ p$ its value trace. Then
\begin{equation}\label{eq:action}
 g'=\gamma\cdot p',\qquad \oint\gamma\cdot dp=0,
\end{equation}
and, along the spherical curve $\widehat q(s)=q(p(s))$,
\begin{equation}\label{eq:tangent-sign}
 B_H(\widehat q',\widehat q')=\frac{p'\cdot\gamma'}{w(p)}.
\end{equation}
We call a trace \emph{positive} when $p'\cdot\gamma'>0$.
Such a curve is noncharacteristic. A collar is an open tubular
neighborhood of a specified curve (or a one-sided neighborhood at an
actual boundary). A germ is a function on a collar, considered up to restriction
to a smaller collar. Retaining a germ means agreement on an open overlap,
so all derivatives agree at the eventual join.

\subsection{Global injectivity, duality, and affine transport}
Local negative determinant will give local inversion. The following
degree argument is what upgrades it to global injectivity.
\par\noindent\textbf{Formalization: proved.} See the precise object and dependency interfaces in the blueprint.\par
\begin{lemma}[Annular degree criterion]\label{lem:degree}
Let $F$ be a smooth map from a compact annulus to $\R^2$. Suppose the
interior Jacobian has a constant nonzero sign, the boundary maps are
homeomorphisms onto disjoint nested Jordan curves, and their induced
winding numbers are opposite. Then $F$ is a diffeomorphism of interiors
onto the intervening planar annulus and a homeomorphism of closures.
The boundary Jacobian may vanish.
\end{lemma}
\begin{proof}
For a target off the boundary images, all preimages are interior and the
degree is defined by compactness. The boundary winding sum has absolute
value one between the curves and is zero in the complementary regions.
All local degrees have the same sign, so there is exactly one preimage
between the curves and none outside. No interior point maps to a boundary
image: local openness there would produce a target on the forbidden side.
Together with the boundary hypotheses this gives a continuous bijection
of compact closures, hence a homeomorphism. Local inversion gives the
smooth interior inverse.
\end{proof}

If $\nabla G$ is a diffeomorphism, its \emph{differential Legendre transform}
is
\[
 G^*(y)=z\cdot y-G(z),\qquad y=\nabla G(z).
\]
If $z=(\nabla G)^{-1}(y)$, differentiation gives
$\nabla G^*(y)=z$ and $D^2G^*(y)=(D^2G(z))^{-1}$.
In particular, $\det D^2G^*(y)=1/\det D^2G(z)<0$ on a saddle domain.
The transform is defined on the gradient image, and its inverse
transform recovers $G$ with the chosen additive constant. This use of
Legendre duality requires a global gradient inverse on the domain in
question; negative determinant alone supplies only a local inverse.

We fix the planar maps and unit vectors
\[
 J(x_1,x_2)=(-x_2,x_1),\qquad S(x_1,x_2)=(x_1,-x_2),
 \qquad A_*=JS,\qquad e_\theta=(\cos\theta,\sin\theta).
\]
Thus $e_\theta'=Je_\theta$ and $SJ=-JS$.
For an invertible affine map $x\mapsto Ax+b$, the correct transport of
a bending is
\begin{equation}\label{eq:affine-transport}
 \widetilde X=AX+b,\qquad \widetilde y=A^{-T}y.
\end{equation}
Derivative pairings are unchanged, so $S_{\widetilde X}\widetilde y=S_Xy$.
With $\widetilde N=A^{-T}N/|A^{-T}N|$, the second fundamental form is
$\II_{\widetilde X}=\II_X/|A^{-T}N|$. Its scalar factor is positive.
An affine change therefore preserves curvature signs and asymptotic
directions, while generally changing the metric and the curvature values.
The formula for $\widetilde y$ preserves the strain equation; it does
not claim that $AX$ and $X$ have the same metric. Affine maps also preserve
tightness, since inverse images of half-spaces are half-spaces.
See \cite{IzmProj} for the broader projective rigidity correspondence.

\section{Identity return and localized bendings}\label{sec:theory}
\label{sec:ruled}
This section isolates the local mechanism. Holonomy here means return of
an asymptotic line field, not Levi--Civita parallel transport. Identity
return is compatible with strictly negative curvature. It supplies a
transverse leaf parameter, and the ruled equations turn that parameter
into the full supported bending space.

\subsection{A transverse coordinate for closed leaves}
\par\noindent\textbf{Formalization: pending.} Deferred from the first-pair objective.\par
\begin{proposition}[Product and first-integral criterion]\label{prop:product-holonomy}
For a smooth nonsingular line foliation near an embedded closed leaf in
an annulus, the following conditions are equivalent after shrinking the
neighborhood: the return germ is the identity; the foliation is a product
$S^1\times J$ by circles; there is a smooth submersion whose connected
fibers are the leaves. Equivalently, a nonvanishing defining one-form has
a nonvanishing integrating factor that makes it exact.
\end{proposition}
\begin{proof}
Start with a coordinate on a short transversal and extend it by transport
along leaves. It is single-valued after one circuit exactly when return
is the identity. It then defines the submersion and the local product.
A product has identity return. The differential of the submersion is a
nonzero multiple of any defining form. Conversely, an exact nonvanishing
defining form is the differential of a transverse coordinate; near the
original compact leaf, its level sets are circles.
\end{proof}
The proposition concerns only the foliation. A general saddle surface
need not have a supported bending for every transverse function. The
ruled model below provides that additional fact.

\subsection{The principal-normal model and its return}
Let $\gamma(s)$ be a unit-speed closed curve of period $L$. Assume its
periodic positively oriented frame $(T,E,n)$ satisfies
\begin{equation}\label{eq:frame}
 T_s=kE,\qquad E_s=-kT+\tau n,\qquad n_s=-\tau E,
 \qquad \tau\ne0.
\end{equation}
The curvature $k$ is signed and may vanish. Here $\tau$ is the
smooth frame coefficient in \eqref{eq:frame}. It agrees with Frenet
torsion where $k\ne0$; the frame equations define it through the
inflection points. Set
\[
 X(s,u)=\gamma(s)+uE(s),\quad \rho=\sqrt{|\tau|},\quad
 \lambda=\tau_s/\tau,\quad D=k_s-k\lambda.
\]
For the normal determined by $X_s\times X_u$, differentiation gives
\begin{align}\label{eq:ruled-II}
 g_X&=\mathcal E\,ds^2+du^2,
 &\mathcal E&=(1-ku)^2+\tau^2u^2,\\
 \II_{ss}&=\frac{u\tau(\lambda+uD)}{\sqrt{\mathcal E}},
 &\II_{su}&=\frac\tau{\sqrt{\mathcal E}},
 &\II_{uu}&=0.\nonumber
\end{align}
In particular $K=-\tau^2/\mathcal E^2<0$. One asymptotic family consists
of the rulings. The other satisfies
\[
 u_s=-\frac u2(\lambda+uD),\qquad
 \frac d{ds}\frac1{\rho u}=\frac D{2\rho}.
\]
The reciprocal coordinate turns the return calculation into a single
period. Put
\[
 I_0=\int_0^L\frac D{2\rho}\,ds.
\]
Periodicity of $\rho$ gives
\begin{equation}\label{eq:return}
 \mathcal R(u)=\frac{u}{1+\rho(0)I_0u}.
\end{equation}
Thus $\mathcal R'(0)=1$ in every case, but the full return is the identity
only if $I_0=0$.

There is also a transverse coordinate regular at the central leaf.
Choose periodic $\omega$ with $\omega_s=D/(2\rho)-I_0/L$. Then
\begin{equation}\label{eq:regular-holonomy-coordinate}
 z=\frac{\rho u}{1-\rho u\omega(s)}
\end{equation}
has $z_u(s,0)=\rho(s)>0$ and changes the characteristic field to
$\partial_s-(I_0/L)z^2\partial_z$. When $I_0=0$, $z$ is a smooth first
integral across $u=0$; on a one-sided band, $v=z^{-1}$ is more convenient.

\begin{remark}[What the period measures]\label{rem:period-holonomy}
The form $(D/(2\rho))\,ds$ has one cohomological obstruction to a
periodic primitive: its integral $I_0$. In this principal-normal class,
$\mathcal R''(0)=-2\rho(0)I_0$, so the quadratic jet detects the whole
return. Along the central curve,
$\II=\left(\begin{smallmatrix}0&\tau\\\tau&0\end{smallmatrix}\right)$
in the orthonormal basis $(T,E)$. The transverse normal curvature is zero,
which explains the automatic multiplier one. This does not say that the
whole band is minimal.
\end{remark}

\subsection{Classification of the supported kernel}
\par\noindent\textbf{Formalization: proved.} See the precise object and dependency interfaces in the blueprint.\par
\begin{theorem}[Supported kernel on a ruled band]\label{thm:ruled}
On a sufficiently thin one-sided band $0<u<b$, a nonzero compactly
supported infinitesimal bending exists exactly when $I_0=0$. In that
case choose periodic $\omega$ with $\omega_s=D/(2\rho)$ and define
\[
 v(s,u)=\frac1{\rho(s)u}-\omega(s),\qquad
 c_b=\max_s\left(\frac1{\rho(s)b}-\omega(s)\right).
\]
Every supported bending has a unique expression
\begin{equation}\label{eq:yF}
 y_F=\frac{\tau u}{2\rho}F'(v)T+
 \left[uF(v)-\frac{1-ku}{2\rho}F'(v)\right]n,
 \qquad F\in C_c^\infty((c_b,\infty)).
\end{equation}
\end{theorem}
\begin{proof}
\emph{Step 1: reduce the strain to transport.}
Write $y=\alpha T+\chi E+\beta n$. The $uu$ equation is $\chi_u=0$;
compact support therefore gives $\chi=0$. The remaining equations are
\[
 (1-ku)\alpha_s+u\tau\beta_s=0,\qquad
 (1-ku)\alpha_u+u\tau\beta_u+k\alpha-\tau\beta=0.
\]
For $p=(1-ku)\alpha+u\tau\beta$, these are equivalent to
\begin{equation}\label{eq:transport-p}
 \alpha=p-\frac u2p_u,\qquad
 \beta=\frac{kp+(1-ku)p_u/2}{\tau},\qquad
 p_s-\frac u2(\lambda+uD)p_u=-uDp.
\end{equation}
Consequently $f=p/(\tau u^2)$ is constant along the nonruling
characteristics.

\emph{Step 2: support forces identity return.}
Extend the field by zero to $u>0$. A characteristic on which $f\ne0$
cannot leave the compact support: transport preserves its nonzero value.
It exists for all $s$ and stays in a compact interval bounded away from
$u=0$. But $(\rho u)^{-1}$ gains $I_0$ on every circuit, which is
incompatible with this boundedness unless $I_0=0$.

\emph{Step 3: reconstruct all supported fields.}
When $I_0=0$, the complete leaves lying in the band are
\begin{equation}\label{eq:closed-leaves}
 u_c(s)=\frac1{\rho(s)(c+\omega(s))},\qquad c>c_b.
\end{equation}
Taking $p=\tau u^2F(v)$ in \eqref{eq:transport-p} gives \eqref{eq:yF}.
A compact set of levels gives a compact subannulus, so every displayed
field is supported. Conversely, if a field is supported in
$u_-\le u\le u_+$ with $0<u_-<u_+<b$, every nonzero level satisfies
\[
 \max_s\left(\frac1{\rho(s)u_+}-\omega(s)\right)
 \le c\le
 \min_s\left(\frac1{\rho(s)u_-}-\omega(s)\right).
\]
The left endpoint is greater than $c_b$ and the right endpoint is finite.
Evaluation of $f$ at $s=0$ therefore gives a smooth compactly supported
function $F$ on $(c_b,\infty)$. Transport proves reconstruction, and every
such level occurs on a complete leaf, proving uniqueness.
\end{proof}
We call a subband \emph{saturated} when it is a union of complete leaves.
Disjoint transverse supports in \eqref{eq:yF} give disjoint surface
supports. This feature, rather than just the dimension of the kernel,
is what will permit exact independent sign choices.


Figure~\ref{fig:bending} gives a geometric reading of
\eqref{eq:yF}. A test function of $v$ selects whole closed leaves,
so the support is an annulus rather than a segment of one characteristic.
Its derivative also enters the field: $y_F$ is a coupled tangential and
normal displacement, with no ruling component. The strain vanishes,
but the quadratic metric correction need not vanish. Consequently the
opposite signs give the common metric described in
\eqref{eq:intro-quadratic}, rather than an isometric path through $X$.
\begin{figure}[htbp]
\centering
\includegraphics[width=\textwidth]{figures/04-bending-map.pdf}
\caption{The bending map $F\mapsto y_F$ in Theorem~\ref{thm:ruled}.
(a) Cutting the band at $s=0$ gives a cylinder diagram with $s=0$ and
$s=L$ identified. Horizontal curves have constant first integral $v$;
vertical curves correspond to rulings at fixed $s$.
(b) A smooth function supported between $v_1$ and $v_2$ selects the gold
subannulus; the field is zero outside it.
(c) At each point the exact field has components
$\langle y_F,T\rangle=\tau uF'(v)/(2\rho)$ and
$\langle y_F,n\rangle=uF(v)-(1-ku)F'(v)/(2\rho)$, while
$\langle y_F,E\rangle=0$.
The displayed arrow illustrates this decomposition; its length and angle
are schematic. The characteristic transport equation, not the arrow,
establishes the bending identity.}
\label{fig:bending}
\end{figure}
\FloatBarrier


\subsection{The return multiplier and the Han--Khuri integral}
For later applications we record what the linear return measures on an
arbitrary saddle surface.
\par\noindent\textbf{Formalization: pending.} Deferred from the first-pair objective.\par
\begin{lemma}[Asymptotic return multiplier]\label{lem:holonomy-integral-new}
Let $\Gamma$ be an embedded closed asymptotic curve on an oriented
surface with $K<0$ near $\Gamma$. Parametrize it by oriented arclength,
write $T$ for its tangent and $Z$ for a unit perpendicular tangent, and set
\[
 k_g=\langle\nabla_TT,Z\rangle,\qquad k_n=\II(Z,Z),\qquad
 \sigma=\operatorname{sign}\II(T,Z).
\]
If $\mu$ is the derivative of the asymptotic return at $\Gamma$, then
\begin{equation}\label{eq:holonomy-integral-new}
 \log\mu=\frac\sigma2\oint_\Gamma k_gk_n|K|^{-1/2}\,ds.
\end{equation}
\end{lemma}
\begin{proof}
Choose Fermi coordinates $(s,t)$ with metric $h^2ds^2+dt^2$ and
$h(s,0)=1$, and write $\II=L\,ds^2+2M\,ds\,dt+N\,dt^2$.
On $t=0$ we have $L=0$, $N=k_n$, and $M=\sigma\sqrt{-K}$.
Codazzi gives $L_t=M_s-k_gN$. If the asymptotic slope is $t_s=f(s,t)$,
then $f(s,0)=0$ and $L+2Mf+Nf^2=0$ implies $f_t=-L_t/(2M)$ on the
leaf. The variational equation yields
\[
 \log\mu=-\frac12\oint\frac{L_t}{M}\,ds
 =-\frac12\oint\frac{M_s}{M}\,ds+
   \frac12\oint\frac{k_gN}{M}\,ds.
\]
The first integral is zero by periodicity. Substitution of $M$ and $N$
in the second gives the assertion.
\end{proof}
The integral in \cite{HK} detects the linear multiplier. Equation
\eqref{eq:return} explains why multiplier one alone does not imply full
identity return.


The identity-return mechanism also produces the visibly broader examples
in Figure~\ref{fig:alternative-bands}. In spherical arclength $r$,
constant spatial speed $a=1$ gives
$I_0=\tfrac12\oint\kappa_r\,dr=0$. Closure still requires the separate
vector moment $\oint aP\,dr=0$. These examples have different normal
tracks from the hemispherical seed used for the torus; they illustrate
the return calculation without supplying its visibility and completion
package.
\begin{figure}[htbp]
\centering
\includegraphics[width=\textwidth]{figures/02-alternative-bands.pdf}
\caption{Three-, four-, and five-lobed constant-speed ruled bands from
the numerical atlas. All axes in each model use equal physical units;
the broad strips come from actual ruling widths, not a transverse
stretch. Gold is the closed spatial center, colored curved traces show
the closed nonruling asymptotic family, and pale straight segments show
the ruling family. Traces are overlaid to make the foliation visible.
With $a=1$, the nonruling equation is $u_r=-\kappa_r u^2/2$, whose
reciprocal gains zero over a period. These are illustrations by immersed
bands: self-intersections are allowed in the saved models. They are not
asserted to be embedded protected cores or completed tight tori.}
\label{fig:alternative-bands}
\end{figure}


\section{Closing a normal loop and cancelling its return period}
\label{sec:structural-balance}
We now construct the frame used in Section~\ref{sec:theory}. The input is
a loop of normals, not a spatial curve. A positive speed determines the
spatial primitive. Closure and return then become two explicit conditions
on that speed, which can be solved separately.

\subsection{The two conditions on speed}
Let $\zeta:\R/\mathcal L\mathbb Z\to\Sph^2$ be a regular spherical loop
parametrized by spherical arclength $r$. Define
\begin{equation}\label{eq:structural-normal-frame}
 P=-\zeta\times\zeta_r,\qquad
 \zeta_{rr}=-\zeta-\kappa P,\qquad P_r=\kappa\zeta_r.
\end{equation}
These formulas fix the sign of its geodesic curvature $\kappa$.
A positive function $a$ determines $\Gamma_r=aP$. If $ds=a\,dr$,
then the frame $(P,\zeta_r,\zeta)$ has $k=\kappa/a$ and $\tau=-1/a$.
Thus
\begin{equation}\label{eq:structural-moments}
 \mathbf M(a)=\oint aP\,dr=0,\qquad
 \mathscr B(a)=\oint\frac{\kappa_r}{\sqrt a}\,dr=0,
 \qquad I_0=\frac12\mathscr B(a).
\end{equation}
The first equation closes the spatial curve; the second cancels the full
return period. Using $-aP$ gives the other torsion sign with the same
conditions.

\subsection{A local speed correction that preserves moments}
The next lemma makes an arbitrarily small positional change while
prescribing a possibly large nonlinear period. Its smallness is in $L^1$
for the speed, rather than in curvature or torsion norms.
\par\noindent\textbf{Formalization: proved.} See the precise object and dependency interfaces in the blueprint.\par
\begin{lemma}[Finite-moment period prescription]\label{lem:finite-moment-balance}
Let $f:S^1\to\R^m$ and $g:S^1\to\R$ be smooth, with $g$ taking both
signs. For every smooth $b>0$, $B_*\in\R$, and $\eta>0$, there is a
smooth $a>0$ such that
\begin{equation}\label{eq:finite-moment-balance}
 \|a-b\|_{L^1}<\eta,\qquad
 \oint af=\oint bf,\qquad \oint g/\sqrt a=B_*.
\end{equation}
If $d=\dim\operatorname{span}f(S^1)$, the correction uses at most $d+1$
arbitrarily short intervals.
\end{lemma}
\begin{proof}
\emph{Choose a slowdown and compensating controls.}
If the period already equals $B_*$, take $a=b$. Otherwise choose a short
interval $K$ where $g$ has the sign of $B_*-\oint g/\sqrt b$.
Choose $d$ disjoint control intervals $O_j$ away from $K$, with smooth
bumps $\psi_j$ whose moments $\oint\psi_jf$ form a basis of
$\operatorname{span}f(S^1)$. One can first choose independent sample
values of $f$, then a center for $K$ in the required sign region away
from those finitely many points, and finally shrink all intervals.
Let $0\le\chi\le1$ be supported in $K$ and equal one on a smaller arc.
Solve
\[
 \sum_{j=1}^d c_j\oint\psi_jf=\oint\chi bf,
\]
and set, for $0\le t<1$,
\[
 a_t=b-t\chi b+t\sum_{j=1}^d c_j\psi_j.
\]
This path preserves every linear moment. The argument also covers $d=0$,
when there are no controls.

\emph{Keep the correction small and prescribe the period.}
Fixing the controls and shortening $K$ gives $c_j=O(|K|)$. Thus $a_t$
is positive on the controls and $\|a_t-b\|_{L^1}<\eta$, uniformly in
$t$. On $K$, positivity follows from $t<1$. On the plateau,
$a_t=(1-t)b$, so the contribution of $g/\sqrt{a_t}$ diverges with the
chosen sign as $t\uparrow1$. All contributions on $K$ have that sign;
those on the complement stay bounded. Continuity therefore attains $B_*$
at some $t_*<1$. Take $a=a_{t_*}$.
\end{proof}

\subsection{The structural existence criterion}
\par\noindent\textbf{Formalization: proved.} See the precise object and dependency interfaces in the blueprint.\par
\begin{theorem}[Normal-loop existence criterion]\label{thm:normal-loop-criterion}
For the loop in \eqref{eq:structural-normal-frame}, the following are
equivalent:
\begin{enumerate}[label=\textup{(\roman*)}]
\item $0$ belongs to the ordinary three-dimensional interior of
$\conv P(S^1)$;
\item a smooth positive speed $b$ satisfies $\mathbf M(b)=0$;
\item a smooth positive speed $a$ satisfies both
$\mathbf M(a)=0$ and $\mathscr B(a)=0$.
\end{enumerate}
This characterizes closed immersed principal-normal identity-return bands
with the prescribed normal loop. For every closing speed $b$, the balanced
speed $a$ can be chosen arbitrarily close to $b$ in $L^1$.
\end{theorem}
\begin{proof}
\emph{From the convex hull to a closing speed.}
Assume (i) and consider
\[
 \Phi(\lambda)=\oint e^{\lambda\cdot P(r)}\,dr,
 \qquad \lambda\in\R^3.
\]
Its Hessian is positive definite because $P$ spans $\R^3$. The interior
condition gives a uniform positive lower bound for
$\max_r v\cdot P(r)$ as $v$ ranges over unit vectors. Uniform continuity
then gives arcs of uniformly positive length on which this scalar product
is bounded below. Hence $\Phi$ is coercive. At its unique minimizer
$\lambda_*$, the positive speed $b=e^{\lambda_*\cdot P}$ has
$\mathbf M(b)=\nabla\Phi(\lambda_*)=0$.

\emph{Why positive closure forces the interior condition.}
Assume (ii). If (i) fails, separation gives a nonzero $v$ with
$v\cdot P\ge0$. Since $b>0$ and the moment is zero, $v\cdot P=0$
everywhere. Differentiation gives $\kappa v\cdot\zeta_r=0$.
On any interval where $\kappa\ne0$, $v$ is perpendicular to both
$P$ and $\zeta_r$, so $\zeta$ is parallel to $v$. This contradicts
$|\zeta_r|=1$. Thus $\kappa=0$ everywhere, making $P$ constant;
a positive speed along a constant unit vector cannot close. This proves (i).

\emph{Cancel the period without changing closure.}
A closing speed also forces $\kappa$ to be nonconstant. Indeed, if
$\kappa=c$, then $C=P-c\zeta$ is constant and $P\cdot C=1$, giving
$\mathbf M(b)\cdot C=\oint b>0$. Therefore $\kappa_r$ takes both signs.
Apply Lemma~\ref{lem:finite-moment-balance} with $f=P$, $g=\kappa_r$,
and $B_*=0$. This gives (iii) with arbitrary $L^1$ control; (iii) implies
(ii) immediately. The frame formulas and \eqref{eq:structural-moments}
then produce the asserted immersed band, whose curvature is negative.
\end{proof}

\par\noindent\textbf{Formalization: proved.} See the precise object and dependency interfaces in the blueprint.\par
\begin{corollary}[Relative upgrade of an embedded seed]\label{cor:relative-holonomy-upgrade}
If a closing speed $b$ has an embedded primitive $\Gamma_b$, the balanced
primitive can be chosen embedded and arbitrarily uniformly close to it.
The two are connected through embedded closed curves with the same oriented
tangent and normal at each parameter value. Any specified regular embedded
linear projection remains embedded. If $\zeta$ is embedded, the Gauss map
of a sufficiently thin resulting band is injective.
\end{corollary}
\begin{proof}
Use the same initial point for all primitives. Their uniform distance is
bounded by the $L^1$ change of speed. On a short parameter arc, some fixed
linear functional is strictly increasing along every positive-speed
primitive of $P$; this excludes nearby collisions. Compact separation for
$\Gamma_b$ excludes all other collisions under a small uniform change.
The argument holds along the whole correcting path and for a specified
regular embedded projection. Tangent and normal directions are fixed.
Finally the ruled Gauss map is locally invertible and injective on the
compact central curve; it is injective on a thin neighborhood.
\end{proof}
An embedded closed asymptotic curve on an oriented $K<0$ surface has
nonzero geodesic torsion \cite[\S2.2]{GR}. Its normal, parametrized by
spherical arclength, gives the frame above with one torsion sign. The
corollary upgrades it to an identity-return ruled band without changing
its knot type or normal self-linking. It does not deform the original
surface isometrically. The global completion still needs the two visible
traces supplied in the next section.

\begin{remark}[The signed action of speed and curvature]\label{rem:holonomy-action}
With $\rho=a^{-1/2}$, the return period is $2I_0=\oint\rho\,d\kappa$.
We call this the return action of the closed planar curve
$r\mapsto(\kappa(r),\rho(r))$. It equals the negative of its oriented
algebraic area for the standard orientation $d\kappa\wedge d\rho$.
A positive speed law $a=A(\kappa)$ cancels
this action automatically, but need not solve the independent closing
equation $\oint A(\kappa)P\,dr=0$. The correction lemma can prescribe
any return period while preserving closure; adding $1$ to the linear
moments also preserves total length. The positional control remains
$C^0$, obtained from $L^1$ control of speed.
\end{remark}

\section{Constructing the embedded protected core}\label{sec:protected-core}
The structural criterion guarantees an immersed identity-return band.
For completion we need more: embedded spatial and normal curves, an
embedded horizontal projection, and two visible planar traces. The seed
below supplies these together. Its long quantitative verification is
deferred to Appendix~\ref{app:seed-estimates}.

\subsection{Visibility and the seed package}\label{sec:corrugated-seed}
For $|z|>R>0$ define the unit vector
\begin{equation}\label{eq:visibility}
 V_R(z)=\frac{Rz+\sqrt{|z|^2-R^2}\,Jz}{|z|^2}.
\end{equation}
Its numerator consists of perpendicular vectors with squared lengths
summing to $|z|^4$.
\begin{definition}\label{def:visible}
Regular closed planar curves $(p,\delta)$ with a common parameter are
\emph{visible at radius $R$} if
\begin{equation}\label{eq:visible-pair}
 |\delta|>R,\qquad p'\cdot V_R(\delta)>0,
 \qquad \delta'\cdot V_R(\delta)>0.
\end{equation}
\end{definition}
These conditions persist under common orientation-preserving
reparametrization. On a compact parameter circle their strict margins
also persist under small $C^1$ changes, or small positional changes with
oriented tangent directions fixed. The connector construction will use
visibility to replace a complicated gradient trace by a circle.

Identify $\R^2$ with $\mathbb C$. For $N\ge2$, set
$\epsilon=N^{-1}$, $T_N=2\pi/N$, and $\xi_N=e^{iT_N}$. Define
\begin{equation}\label{eq:corrugated-beta}
 \beta_N(t)=(1+\epsilon\sin Nt)
          \exp\bigl(i[t+2\epsilon\sin Nt]\bigr).
\end{equation}
Choose $k_N$ as the unique root of
\begin{equation}\label{eq:corrugated-root}
 \int_0^{2\pi}e^{-k\cos x}(1+\epsilon\sin x)^2(1+2\cos x)\,dx=0,
\end{equation}
and put
\[
 Z_N=\frac1{2\pi}\int_0^{2\pi}e^{-k_N\cos x}\,dx,
 \qquad m_N(t)=Z_N^{-1}e^{-k_N\cos Nt}.
\]
The partner curve is defined by integrating a positive multiple of the
same tangent, with an initial point chosen for rotational covariance:
\begin{equation}\label{eq:corrugated-delta}
 I_N=\int_0^{T_N}m_N\beta_N'\,dt,
 \qquad
 \delta_N(t)=\frac{I_N}{\xi_N-1}+\int_0^t m_N(s)\beta_N'(s)\,ds.
\end{equation}

\par\noindent\textbf{Formalization: proved.} See the precise object and dependency interfaces in the blueprint.\par
\begin{lemma}[Corrugated analytic seed]\label{lem:seed}
For every integer $N\ge10^4$, $\beta_N$ and $\delta_N$ are real analytic
positively oriented embedded curves enclosing the origin. If $Q_N$ is
rotation through $2\pi/N$, then
\begin{equation}\label{eq:seed-data}
 \beta_N(t+T_N)=Q_N\beta_N(t),\qquad
 \delta_N(t+T_N)=Q_N\delta_N(t),\qquad
 \delta_N'=m_N\beta_N',\quad m_N>0,
\end{equation}
and
\begin{equation}\label{eq:zero-action}
 \int_0^{2\pi}\det(\beta_N,\delta_N')\,dt=0.
\end{equation}
The pair $(\beta_N,\delta_N)$ is visible at $R_+=1/4$, and
\[
 (\widetilde\beta_N,\widetilde\delta_N)(t)
 =\bigl(A_*\delta_N(-t),A_*\beta_N(-t)\bigr)
\]
is visible at $R_-=4/5$. The spherical lift
\[
 \zeta_N(t)=\frac{(\beta_N(t),1)}{\sqrt{1+|\beta_N(t)|^2}}
\]
is embedded in the northern hemisphere and its geodesic curvature takes
both signs.
\end{lemma}
\begin{proof}
Appendix~\ref{app:seed-estimates} verifies the root, covariance, action,
embedding, visibility, and curvature assertions in that order.
\end{proof}

\begin{remark}[The role of corrugation]\label{rem:corrugation-mechanism}
The curves satisfy $\beta_N=e^{it}+O(N^{-1})$ and
$\delta_N=(i/2)e^{it}+O(N^{-1})$ in position, while their common unit
tangent has order-one oscillations. This allows signed angular backtracking
and zero mixed action without losing nearly circular positional geometry.
The latter makes both visibility estimates uniform. The exact threshold
$10^4$ is only a convenient sufficient bound.
\end{remark}


The three views in Figure~\ref{fig:seeds} separate the positional and
tangent effects of corrugation. A modest value of $N$ makes the global
shape legible, whereas a uniformly enlarged single cell at $N=10^4$
shows the small-scale motion within the sufficient analytic range.
The partner is computed from a positive multiple of the same tangent;
its closed primitive is a different curve, not a scaled copy of $\beta_N$.
\begin{figure}[htbp]
\centering
\includegraphics[width=\textwidth]{figures/03-seeds.pdf}
\caption{The seed formulas \eqref{eq:corrugated-beta}--\eqref{eq:corrugated-delta}.
(a) Blue is $\beta_N$ and teal is $\delta_N$ at $N=16$.
(b) The corresponding northern spherical lift $\zeta_N$.
These two panels make corrugation visible; $N=16$ is below the sufficient
range of Lemma~\ref{lem:seed}, so they are illustrations of the formulas.
(c) One full cell at $N=10^4$: each curve is translated to its own initial
point, and both coordinate differences are multiplied by $N$.
This is an isotropic magnification, not a common placement of the two
curves. Arrows show the parameter direction. The oriented tangents agree
at corresponding parameters because $\delta_N'=m_N\beta_N'$ with
$m_N>0$. The root defining $m_N$ cancels the mixed action, while the
covariant initial point closes the partner. The seed estimates are
proved in Appendix~\ref{app:seed-estimates}.}
\label{fig:seeds}
\end{figure}


\subsection{From planar traces to a closed framed curve}\label{sec:model-curve}
Fix one seed and suppress its subscript $N$. Let $Q$ denote rotation by
$2\pi/N$, extended to $\R^3$ by fixing the third coordinate. The
coordinate transitions used in this section are summarized here:
\[
\begin{array}{c|l}
 t & \text{original seed parameter, of period }2\pi\\
 r & \text{spherical arclength of }\zeta,\text{ of period }N\ell\\
 a(r) & \text{positive spatial speed: }\Gamma_r=aP\\
 s & \text{spatial arclength: }ds=a\,dr,\text{ of period }L\\
 u & \text{distance along the ruling }E=\zeta_r\\
 v & \text{transverse first integral labelling closed leaves}
\end{array}
\]
Reparametrize $\zeta$ by $r$; covariance gives
$\zeta(r+\ell)=Q\zeta(r)$. Its frame is
\begin{equation}\label{eq:spherical-frame}
 P=-\zeta\times\zeta_r,\qquad
 \zeta_{rr}=-\zeta-\kappa P,\qquad P_r=\kappa\zeta_r.
\end{equation}
Both $\kappa$ and $|\beta|$ are $\ell$-periodic. With
$w=\sqrt{1+|\beta|^2}$, the planar expression is
\begin{equation}\label{eq:P-formula}
 P=\frac{(-J\beta_r,-\det(\beta,\beta_r))}{w^2}.
\end{equation}
Define the spatial lift by
\begin{equation}\label{eq:spatial-lift}
 \Gamma_h=-J\delta,\qquad
 (\Gamma_3)_t=-\det(\beta,\delta_t),
\end{equation}
where $h$ denotes horizontal components. The vertical integrand is
cell-periodic. Its total integral is zero by \eqref{eq:zero-action}, so
each of the $N$ equal cell integrals is zero. Thus the vertical primitive
is periodic on every cell. Fix its additive constant. Since
$\delta_r=m\beta_r$, \eqref{eq:P-formula} gives
\begin{equation}\label{eq:initial-speed}
 \Gamma_r=a_0P,\qquad a_0=mw^2>0.
\end{equation}
Its horizontal projection is embedded, so $\Gamma$ is embedded.

For any positive $\ell$-periodic speed $a$, closure reduces to the scalar
moment
\begin{equation}\label{eq:closing-moment}
 M(a)=\int_0^\ell aP_3\,dr=0.
\end{equation}
Indeed, the horizontal cell increments are successive rotations and
sum to zero; the vertical increments are all equal. Choose the horizontal
initial point by
\begin{equation}\label{eq:covariant-start}
 (Q-I)\Gamma_h(0)=\int_0^\ell aP_h\,dr.
\end{equation}
Then $\Gamma(r+\ell)=Q\Gamma(r)$. The new planar trace
$\delta_a=J\Gamma_h$ satisfies $(\delta_a)_r=(a/w^2)\beta_r$:
changing speed preserves its oriented tangent direction. The initial
speed recovers $\delta$ and satisfies $M(a_0)=0$.

The return balance over one cell is
\begin{equation}\label{eq:balance-moment}
 \mathscr B_\ell(a)=\int_0^\ell\frac{\kappa_r}{\sqrt a}\,dr.
\end{equation}
We use $\mathscr B_\ell$ here to distinguish this scalar from the support
tensor $B_H$.

\par\noindent\textbf{Formalization: proved.} See the precise object and dependency interfaces in the blueprint.\par
\begin{proposition}[One-slowdown holonomy correction]\label{prop:one-slowdown}
If $b>0$ is $\ell$-periodic with $M(b)=0$, $\kappa$ is nonconstant,
and $P_3\not\equiv0$, then for every $\eta>0$ a positive smooth
$\ell$-periodic $a$ satisfies
\[
 \|a-b\|_{L^1(0,\ell)}<\eta,\qquad M(a)=0,
 \qquad\mathscr B_\ell(a)=0.
\]
The change is supported in at most two arbitrarily short intervals.
\end{proposition}
\begin{proof}
Apply Lemma~\ref{lem:finite-moment-balance} on the cell circle with
$f=P_3$, $g=\kappa_r$, and $B_*=0$. The moment rank is one, and the
periodic nonconstant $\kappa$ has derivative of both signs. One slowdown
and one compensating control suffice.
\end{proof}

\par\noindent\textbf{Formalization: proved.} See the precise object and dependency interfaces in the blueprint.\par
\begin{theorem}[Balanced embedded speed]\label{thm:spike}\label{lem:speed}
For every $\eta_0>0$, the seed admits a positive smooth $\ell$-periodic
speed $a$ with
\begin{equation}\label{eq:balanced-speed}
 \|a-a_0\|_{L^1(0,\ell)}<\eta_0,\qquad
 M(a)=0,\qquad\mathscr B_\ell(a)=0.
\end{equation}
For small enough $\eta_0$, its spatial primitive and horizontal
projection are embedded, and both visibility conditions persist.
\end{theorem}
\begin{proof}
The seed has nonconstant $\kappa$. Also $P_3\not\equiv0$: otherwise
$\det(\beta,\beta_r)=0$ would prevent $\beta$ from winding around zero.
The proposition therefore gives the speed. For fixed $N$, the initial
point in \eqref{eq:covariant-start} varies continuously with the cell
integral, and integration of $aP$ gives uniform closeness of primitives.
On short parameter arcs, a fixed linear functional is monotone along all
positive-speed primitives; on the remaining pairs, the original embedded
curve has positive separation. This proves embedding for the spatial
curve and its horizontal projection. Visibility, after division by positive
speeds, depends on positions and the fixed oriented tangent directions.
Its strict margins therefore survive the uniform positional change.
The seed parameter $N$ is fixed before choosing the speed correction.
The band width is chosen only after the final positive speed $a$ is
fixed; the $L^1$ estimate does not give uniform bounds for its curvature
or torsion.
\end{proof}

\begin{remark}[Why a constant speed does not suffice]\label{rem:constant-speed-obstruction}
A constant $a$ cancels the return period, but cannot close this
hemispherical normal loop. If $D$ is the spherical disk bounded by
$\zeta$, then
\[
 \oint\Gamma_r\,dr=-a\oint\zeta\times d\zeta
 =-2a\int_D\zeta\,dA.
\]
The vertical component of the last integral is positive. Thus speed
redistribution solves a real closure obstruction.
\end{remark}

\subsection{The balanced band and its support tensor}
Fix the final speed and set $ds=a\,dr$, with
$L=\int_0^{N\ell}a\,dr$. In physical arclength the frame is
$T=P$, $E=\zeta_r$, $n=\zeta$, and
\begin{equation}\label{eq:frame-speed}
 k=\frac\kappa a,\qquad \tau=-\frac1a,\qquad
 \lambda=-\frac{a_r}{a^2},\qquad D=\frac{\kappa_r}{a^2}.
\end{equation}
These identities follow from $\partial_s=a^{-1}\partial_r$; in particular
the subscripts defining $\lambda$ and $D$ are $s$ derivatives. Hence
\[
 I_0=\int_0^{N\ell}\frac{\kappa_r}{2\sqrt a}\,dr
 =\frac N2\mathscr B_\ell(a)=0.
\]
Theorem~\ref{thm:ruled} supplies the localized kernel.

The ruled Gauss map is locally invertible because $K<0$, and its central
restriction is the embedded curve $\zeta$. It is therefore injective on
a thin band. Shrink the band so that its normals remain in the northern
hemisphere. Horizontal projection is then locally invertible and is
injective on the central curve, so it too is injective on a thin band.
Fix a width satisfying both requirements.

\par\noindent\textbf{Formalization: proved.} See the precise object and dependency interfaces in the blueprint.\par
\begin{proposition}[Central support tensor]\label{prop:central-support}
In Gauss coordinates, the support tensor along the central curve is
\begin{equation}\label{eq:central-B}
 B_H=\begin{pmatrix}0&a\\a&0\end{pmatrix}
\end{equation}
in the spherical orthonormal basis $(\zeta_r,P)$. Its mixed entry is
strictly positive.
\end{proposition}
\begin{proof}
The central derivative is $\Gamma_r=aP$, so $B_H\zeta_r=aP$.
At $u=0$, the ruled normal has $N_u=P/a$, while $X_u=\zeta_r$.
Consequently $B_HP=a\zeta_r$, giving the second column.
\end{proof}

\section{Positive exits without changing the bending region}\label{sec:exits}
The ruled core has closed characteristics, whereas the completion needs
positive noncharacteristic seams. We create those seams in two collars
where all retained bendings vanish. The core itself stays fixed.

\subsection{Changing return and finding a positive curve}
\par\noindent\textbf{Formalization: pending.} See the precise object and dependency interfaces in the blueprint.\par
\begin{lemma}[Exit perturbation]\label{lem:exit}
Let a saddle support in the northern hemisphere have a positively oriented
embedded closed asymptotic curve enclosing the north pole. Suppose its
return germ is the identity and the mixed support entry, evaluated on
unit tangent and outward conormal, is positive. Inside any prescribed
collar, an arbitrarily small $C^2$ perturbation fixing the first jet on
the curve can make return attracting or repelling. The attracting choice
gives arbitrarily close positive closed curves on the exterior side;
the repelling choice gives them on the interior side.
\end{lemma}
\begin{proof}
\emph{Step 1: express the multiplier through the normal second jet.}
Parametrize the curve $q_0(r)$ by spherical arclength, of period
$\mathcal L$, and let $P=-q_0\times(q_0)_r$ be its outward conormal.
The Fermi map and metric are
\[
 q(r,t)=\cos t\,q_0(r)+\sin t\,P(r),\qquad
 h^2dr^2+dt^2,\qquad h=\cos t+\kappa(r)\sin t.
\]
Here $t>0$ is exterior and $(q_0)_{rr}=-q_0-\kappa P$.
Write $B_H=L\,dr^2+2M\,dr\,dt+N\,dt^2$, with
$(L,M,N)|_{t=0}=(0,a,c)$ and $a>0$. The coordinate expressions are
\[
 L=H_{rr}-\frac{h_r}{h}H_r+hh_tH_t+Hh^2,
 \quad M=H_{rt}-\frac{h_t}{h}H_r,\quad N=H_{tt}+H.
\]
On the seam they give
\begin{equation}\label{eq:seam-derivative}
 L_t=a_r+\kappa c.
\end{equation}
For clarity, $a_r=H_{rrt}-\kappa_rH_r-\kappa H_{rr}$ there;
subtracting this from $L_t$ leaves
$\kappa(H_{rr}+\kappa H_t+H)+\kappa(H_{tt}+H)=\kappa c$.
The asymptotic slope vanishing on the seam is
\[
 t_r=f(r,t)=-\frac{L}{M+\sqrt{M^2-LN}},\qquad
 f_t(r,0)=-\frac{a_r+\kappa c}{2a}.
\]
Integrating its variational equation gives
\begin{equation}\label{eq:multiplier}
 \mu=\exp\left(-\frac12\int_0^{\mathcal L}\frac{\kappa c}{a}\,dr\right),
\end{equation}
because the period of $a_r/a$ is zero.

\emph{Step 2: change that jet while retaining the curve.}
Choose a cutoff $\chi(t)$ supported in the prescribed collar and equal
one near zero. Put
\[
 H_\epsilon=H+\frac\epsilon2a(r)\kappa(r)t^2\chi(t).
\]
The first jet is unchanged; $L=0$ and $M=a$ remain, while $c$ changes
to $c+\epsilon a\kappa$. Since the original multiplier is one,
\begin{equation}\label{eq:perturbed-multiplier}
 \mu_\epsilon=\exp\left(-\frac\epsilon2
                  \int_0^{\mathcal L}\kappa^2\,dr\right).
\end{equation}
The integral is positive: otherwise the curve would be a great circle,
which cannot lie in an open hemisphere. Thus $\epsilon>0$ gives
attraction and $\epsilon<0$ repulsion. For fixed collar and cutoff, small
$\epsilon$ preserves the saddle determinant by $C^2$ continuity.

\emph{Step 3: produce a positive closed graph on the chosen side.}
Fix such a nonzero $\epsilon$, set $b(r)=f_t(r,0)$ for the perturbed
support, and let $\bar b=(\log\mu_\epsilon)/\mathcal L$.
The zero-mean equation $v_r=(b-\bar b)v$ has a positive periodic solution.
For the graph $t=\delta v(r)$, uniform Taylor expansion gives
\[
 B_{H_\epsilon}(\partial_r+t_r\partial_t,
                       \partial_r+t_r\partial_t)
 =-2ab\delta v+2a\delta v_r+O(\delta^2)
 =-2a\bar b\delta v+O(\delta^2).
\]
This is positive for small $\delta>0$ in the attracting case and small
$\delta<0$ in the repelling case. The graphs are embedded closed curves
in the collar and converge smoothly to the seam.
\end{proof}

\subsection{Fixing the protected annulus and both exits}
Choose $I=(c_-,c_+)$ with compact closure in the complete-leaf interval
of Theorem~\ref{thm:ruled}. Let $V_0$ be the compact annulus swept out
by the levels $c\in[c_-,c_+]$. All fields from $C_c^\infty(I)$ are
supported in $V_0$. Choose levels
$c_b<c_1<c_-<c_+<c_2$. Their leaves bound a larger annulus, and their
disjoint collars can be chosen separated from a neighborhood of $V_0$.
Every retained bending vanishes on those collars.

In the Gauss domain, the leaves are nested positively oriented curves
around the north pole. Their mixed entries remain positive by
\eqref{eq:central-B} and thinness. Apply the exit lemma with attraction
on the outer leaf and repulsion on the inner leaf. Choose the perturbations
small enough to preserve embedding of the support immersion and its
horizontal projection; then choose the positive graphs sufficiently close
to the perturbed seams.

In planar coordinates denote the resulting source traces by $p_-,p_+$,
and their gradient traces by $\gamma_-,\gamma_+$. They satisfy
\begin{equation}\label{eq:positive-traces}
 p_\pm'\cdot\gamma_\pm'>0.
\end{equation}
Both source and gradient maps are embeddings on a neighborhood of the
annulus between them. The individual traces are positively oriented and
enclose zero. The outer pair $(p_+,J\gamma_+)$ is visible at $R_+$;
the inner pair needed after duality is
\begin{equation}\label{eq:reflected-visible}
 \bigl(S\gamma_-(-s),JS p_-(-s)\bigr),
\end{equation}
visible at $R_-$. These follow by continuity from the seed inequalities:
the central source and gradient are $\beta$ and $-J\delta$, and
$SJ=-JS$ identifies the reflected central pair with the one in
Lemma~\ref{lem:seed}.

The source order is $p_-$ inside $p_+$. The gradient order is reversed:
$\gamma_+$ lies inside $\gamma_-$. Indeed, the gradient is an
orientation-reversing diffeomorphism, while both gradient boundary
curves are parametrized counterclockwise; the induced boundary
orientations determine which one is inner. Thus the minus sign in a
trace name refers to its source boundary, not to its gradient position. We fix all these choices in order: final speed, thin
band, protected levels and collars, small exit perturbations, then close
positive exit graphs. Every required margin is strict on a compact set.
From this point onward, an open neighborhood of $V_0$ is kept unchanged.

\section{Support completion by dual radial filling}\label{sec:completion}
The protected band has two positive noncharacteristic exit curves.
The only nonradial extension tool we retain is a visible connector: it
makes an outgoing \emph{gradient} trace circular. A differential Legendre
transform then makes that circle a \emph{source} seam, where a large
radial quadratic term controls the saddle Hessian. On one end we continue
to a polar circle; on the other we attach a square-root Legendre profile
which becomes a reflectable parabolic neck. The same radial lemma treats
both ends, but the roles of source and gradient are exchanged at each
Legendre transform. We record those domains explicitly in the final
completion proof. All changes to the support occur outside the
protected core. A single ambient normalization, applied also to the
retained fields, is recorded in Section~\ref{sec:model-saddle}.


Figure~\ref{fig:completion} follows the completion in two complementary
views. The physical windows retain the actual band and its attachments;
the domain diagram records which circular trace is a gradient and which
is a source. At the final common physical scale the band and neck are
microscopic, so the global cutaway serves to show the assembly, while
the local windows show the attachment geometry.



\subsection{Relative smoothing at a positive seam}
\par\noindent\textbf{Formalization: proved.} See the precise object and dependency interfaces in the blueprint.\par
\begin{lemma}[Relative saddle smoothing]\label{lem:smoothing}
Suppose two smooth planar potentials have matching values and gradients
on a compact embedded closed seam. If their Hessians are indefinite and
their common tangential quadratic form is positive, they admit a smooth
saddle interpolation equal to the old pieces outside any prescribed seam
neighborhood and arbitrarily close to the piecewise potential in $C^1$.
\end{lemma}
\begin{proof}
\emph{The allowed second jets form a convex interval.}
Use tubular coordinates $(s,t)$ with seam $t=0$ and denote the piecewise
function by $G_{\rm pw}$. Matching first jets makes it $C^1$ with locally
Lipschitz first derivatives. Differentiating the common gradient along
the seam shows that the tangential and mixed Hessian entries agree.
In a tangent-normal orthonormal frame the limiting Hessians are
\[
 H_\pm(s)=\begin{pmatrix}a(s)&b(s)\\b(s)&c_\pm(s)\end{pmatrix},
 \qquad a>0,\quad c_\pm<b^2/a.
\]
There is a uniform strict margin, and every convex combination remains
indefinite because only the normal entry varies.

\emph{Mollify and retain the old germs.}
Extend the coordinate cylinder periodically in $s$ and convolve with a
symmetric nonnegative mollifier of radius $\delta$. Since the gradient
has no jump, the distributional Hessian has no seam measure. Close to
the seam, the mollified Hessian is a weighted average of nearby one-sided
Hessians, uniformly $o(1)$ from the displayed convex segments. Conversion
from coordinate Hessians to Euclidean Hessians adds only $o(1)$ errors,
since the gradients converge uniformly. Thus the determinant stays negative.
Away from the seam, ordinary $C^2$ convergence gives the same conclusion.

Patch the mollified function to $G_{\rm pw}$ with a cutoff equal one for
$|t|\le h$ and zero for $|t|\ge2h$, where $h=\sqrt\delta$.
On the transition strips the mollifier stays on one smooth side.
Its value and gradient errors are $O(\delta^2)$ and $O(\delta)$.
The additional Hessian terms are therefore
$O(\delta^2h^{-2}+\delta h^{-1})=o(1)$. The saddle margin survives,
the $C^1$ error tends to zero, and the old pieces are unchanged for
$|t|\ge2h$. Taking $\delta$ small places all changes inside the prescribed
neighborhood.
\end{proof}
When using this lemma we cut the retained boundary inside an unchanged
incoming subcollar. Local smoothing gives no automatic global injectivity;
we verify it from the unchanged boundary images and the degree criterion.

\subsection{A connector with circular outgoing gradient}
The design is to draw a family of straight source segments from $p$,
choose their lengths so the terminal gradient is a circle, and then use
degree to rule out overlaps in both source and gradient.
\par\noindent\textbf{Formalization: pending.} See the precise object and dependency interfaces in the blueprint.\par
\begin{lemma}[From a visible trace to a circular gradient]\label{lem:connector}
Let $p,\delta:\R/T\mathbb Z\to\R^2$ be positively oriented embedded
curves enclosing zero. Put $\gamma=-J\delta$ and assume
\begin{equation}\label{eq:connector-data}
 p'\cdot\gamma'>0,\qquad g'=\gamma\cdot p'
\end{equation}
for a periodic $g$. If $(p,\delta)$ is visible at $R>0$, an exterior
source annulus carries a saddle potential with incoming value $g$ and
gradient $\gamma$, and outgoing gradient a positively parametrized
circle of radius $R$. Both boundaries are positive for the Hessian.
The source and gradient maps are diffeomorphisms of interiors onto their
respective planar annuli. If the terminal gradient is written
$Re_\theta$, its terminal source trace $P_\mathrm{out}(\theta)$ also
satisfies $P_\mathrm{out}\cdot e_\theta>0$ for a sufficiently small
connector parameter. If an incoming saddle germ supplies the data,
relative smoothing retains a smaller incoming collar and the outgoing germ.
\end{lemma}
\begin{proof}
\emph{Step 1: choose the segment directions and lengths.}
Set $\kappa_0=\sqrt{|\delta|^2-R^2}$ and choose an angular lift with
\begin{equation}\label{eq:visibility-angle}
 e_\varphi=V_R(\delta),\qquad
 \delta=Re_\varphi-\kappa_0Je_\varphi.
\end{equation}
Differentiation gives $\varphi'=\delta'\cdot e_\varphi/\kappa_0>0$.
The argument of $\delta$ differs from $\varphi$ by the periodic angle
$-\arctan(\kappa_0/R)$, so $\varphi(s+T)=\varphi(s)+2\pi$.
Choose a small constant $\eta>0$, put $\psi=\varphi-\eta$ and
$w=\delta-Re_\psi$, and define
\begin{equation}\label{eq:connector-scalars}
 E_0=\delta\cdot e_\psi-R,\quad
 A_0=\det(w,p'),\quad B_0=\det(w,\delta'),\quad
 C_0=E_0R\psi',\quad \ell=A_0/C_0.
\end{equation}
All four scalars are positive for small $\eta$, uniformly in $s$:
$E_0=\kappa_0\sin\eta-R(1-\cos\eta)$,
$A_0\to\kappa_0p'\cdot e_\varphi>0$, and
$B_0\to\kappa_0\delta'\cdot e_\varphi>0$.

\emph{Step 2: verify local potential and saddle determinant.}
On $0\le u\le\ell(s)$ put
\begin{equation}\label{eq:connector-map}
 P(s,u)=p+uw,\qquad \Delta=A_0+u(B_0-C_0),\qquad
 \widehat G(s,u)=g+u\gamma\cdot w.
\end{equation}
At the endpoints $\Delta=A_0$ and $\Delta=\ell B_0$, hence $\Delta>0$.
Since $\det(w,w')=B_0-C_0$,
\begin{equation}\label{eq:connector-jacobian}
 \det(P_s,P_u)=-\Delta<0.
\end{equation}
In a local inverse chart $G=\widehat G\circ P^{-1}$ has gradient
\begin{equation}\label{eq:connector-gradient}
 Y=\gamma+\frac{uB_0}{\Delta}Jw,
 \qquad Y_u=\frac{A_0B_0}{\Delta^2}Jw.
\end{equation}
Indeed, $Y\cdot P_u=\widehat G_u$; also $Jw\cdot P_s=\Delta$ and
$\gamma'\cdot w=B_0$, so $Y\cdot P_s=\widehat G_s$.
The pulled-back Hessian has $uu$ entry zero and mixed entry
$A_0B_0/\Delta$. Dividing its determinant by the square of the source
Jacobian gives
\begin{equation}\label{eq:connector-det}
 \det D^2G=-\frac{A_0^2B_0^2}{\Delta^4}<0.
\end{equation}
At $u=\ell$ the fraction in \eqref{eq:connector-gradient} is one, hence
\begin{equation}\label{eq:connector-terminal}
 Y(s,\ell(s))=-RJ e_{\psi(s)}.
\end{equation}

\emph{Step 3: prove global source injectivity.}
As $\eta\downarrow0$, convergence in each fixed smooth norm gives
\begin{equation}\label{eq:connector-limit}
 \eta\ell\longrightarrow z_0:=\frac{p'\cdot e_\varphi}{R\varphi'}>0,
 \qquad \eta P(s,\ell(s))\longrightarrow-z_0\kappa_0Je_\varphi.
\end{equation}
For $\theta=\psi-\pi/2$, the same limit gives
$\eta P(s,\ell(s))\cdot e_\theta\longrightarrow
z_0\kappa_0>0$. This is the radial-derivative condition needed after
Legendre duality. The limit is a radial graph with positive radius and
strictly increasing polar angle $\varphi-\pi/2$. Its radius and angular derivative have
positive minima. For small $\eta$ the terminal source curve is therefore
embedded, positively oriented, and star-shaped, and its unscaled radius
tends uniformly to infinity. It encloses $p$. The induced orientations
of the two boundary curves are opposite. Lemma~\ref{lem:degree}, together
with \eqref{eq:connector-jacobian}, gives global injectivity of $P$ and a
single well-defined potential on its source annulus.

\emph{Step 4: check the seams and global gradient injectivity.}
Differentiating \eqref{eq:connector-limit} and pairing with $e_\psi$ yields
\[
 \eta\frac d{ds}P(s,\ell(s))\cdot e_\psi
 \longrightarrow z_0\kappa_0\varphi'>0.
\]
The outgoing gradient derivative is $R\psi'e_\psi$, so its tangential
Hessian is positive. The incoming one is positive by hypothesis.
The incoming gradient curve $\gamma=-J\delta$ encloses the disk of
radius $R$: it is embedded, surrounds zero, and stays outside that disk.
Thus the gradient boundary curves are disjoint and nested. Equation
\eqref{eq:connector-det} and the degree criterion give global gradient
injectivity. Finally the incoming first jet is exactly $(g,\gamma)$.
Relative saddle smoothing attaches any incoming germ while retaining
both an unchanged incoming subcollar and the outgoing germ. The boundary
images still provide the same degree argument on the joined annulus.
\end{proof}


\subsection{Quadratic filling of a circular source seam}
The next lemma replaces the separate straightening, logarithmic, and cap
constructions. Its proof uses two \emph{linear} inequalities within the
nonconvex saddle condition. For a function in polar coordinates,
\begin{equation}\label{eq:polar-hessian}
 D^2 f\big|_{(e_\theta,Je_\theta)}=
 \begin{pmatrix}
 f_{rr}& f_{r\theta}/r-f_\theta/r^2\\
 f_{r\theta}/r-f_\theta/r^2&(f_{\theta\theta}+rf_r)/r^2
 \end{pmatrix}.
\end{equation}
Thus $f_{rr}<0$ and $f_{\theta\theta}+rf_r>0$ imply
$\det D^2f<0$, regardless of the mixed derivative.

\par\noindent\textbf{Formalization: proved.} See the precise object and dependency interfaces in the blueprint.\par
\begin{lemma}[Quadratic filling at a circular seam]
\label{lem:quadratic-filling}
Let $f$ be a smooth saddle potential on an exterior collar of
$\{|y|=R\}$, with boundary traces
\[
 h(\theta)=f(Re_\theta),\qquad b(\theta)=f_r(Re_\theta).
\]
Assume
\begin{equation}\label{eq:radial-fill-data}
 b(\theta)>0,\qquad h''(\theta)+Rb(\theta)>0.
\end{equation}
There is a smooth saddle extension into $0<|y|<R$, agreeing with $f$
outside any prescribed seam neighborhood, and equal near the puncture to
\begin{equation}\label{eq:quadratic-germ}
 f(y)=c|y|-\frac M2|y|^2+d,\qquad c=MR,
\end{equation}
where $M>0$ can be arbitrarily large. If additionally the trace
$\nabla f(Re_\theta)$ is a positively oriented Jordan curve enclosing
zero, the extension can be chosen with globally injective gradient on
the attached annulus, and the radial inner-gradient circle can be made
to enclose that Jordan curve.
\end{lemma}
\begin{proof}
Choose a smooth function $\chi:[0,R]\to[0,1]$ equal to zero for
$r\le R/2$ and one for $r\ge3R/4$. On $0<r\le R$ put
\begin{equation}\label{eq:quadratic-filling}
 f_M(r,\theta)=-\frac M2(r-R)^2
   +\chi(r)\bigl[h(\theta)+(r-R)b(\theta)\bigr].
\end{equation}
At $r=R$ this has exactly the prescribed value and first derivatives.
All the angular terms and their derivatives are bounded independently
of $M$, while
\begin{equation}\label{eq:quadratic-radial-signs}
 (f_M)_{rr}=-M+O(1),\qquad
 (f_M)_{\theta\theta}+r(f_M)_r
     =Mr(R-r)+O(1).
\end{equation}
At $r=R$, the full second expression is $h''+Rb>0$ by
\eqref{eq:radial-fill-data}. By continuity and the nonnegative
leading term, it remains positive in a fixed narrow collar below
$R$ for large $M$. Off that collar the leading term
$Mr(R-r)$ is positive; on $r<R/2$ the angular error vanishes altogether.
Hence both strict inequalities preceding the lemma hold throughout the
punctured disk. Near $r=0$ formula \eqref{eq:quadratic-filling} is
exactly \eqref{eq:quadratic-germ} with $d=-MR^2/2$.

The two potentials have a common first jet on $r=R$, indefinite
Hessians, and positive common tangential Hessian
$(h''+Rb)/R^2$. Lemma~\ref{lem:smoothing} gives a smooth
interpolation confined to an arbitrarily small seam collar, retaining
both the exterior original germ and the radial inner germ. First fix
$M$, then take that collar sufficiently small. The smoothing preserves
negative Hessian determinant and positive tangential quadratic form;
it need not preserve the sufficient inequality $(f_M)_{rr}<0$ on the
incoming side of the seam.

For the injectivity assertion, write $P(\theta)=\nabla f(Re_\theta)$.
The positive Jordan curve $P$ persists on a nearby unchanged exterior
circle. Take a fixed small inner circle $r=r_0<R/2$. There
\[
 \nabla f_M(r_0e_\theta)=(MR-Mr_0)e_\theta.
\]
For large $M$ this circle surrounds $P$. Between the two boundaries
$\det D^2f<0$, the boundary maps are embedded and nested, and their
induced winding numbers are opposite. The annular degree criterion,
Lemma~\ref{lem:degree}, gives global gradient injectivity, also after
relative smoothing. The same argument applies to any shorter compact
annulus within the extension.
\end{proof}

\subsection{A radial neck adapter}
The quadratic germ can be ended at the puncture, or it can be joined
radially to a square-root germ. The latter has a particularly simple
inverse Legendre transform.

\par\noindent\textbf{Formalization: proved.} See the precise object and dependency interfaces in the blueprint.\par
\begin{lemma}[Square-root neck adapter]\label{lem:neck-adapter}
Suppose a radial saddle potential $f=f(r)$ has $f'>0$ and $f''<0$
near a circle $r=r_j>0$. Choose any $0<a<r_j$ and set
\begin{equation}\label{eq:neck-adapter-constants}
 \ell=r_j-a,\quad B=f'(r_j)^2\ell,\quad
 C=f(r_j)-2f'(r_j)\ell.
\end{equation}
Then
\begin{equation}\label{eq:sqrt-neck}
 f_{\rm neck}(r)=C+2\sqrt{B(r-a)},\qquad a<r\le r_j,
\end{equation}
matches the value and radial derivative of $f$ at $r_j$, and satisfies
$f_{\rm neck}'>0$, $f_{\rm neck}''<0$. Relative radial smoothing gives
a smooth saddle continuation, unchanged near its original incoming
collar and exactly \eqref{eq:sqrt-neck} near $r=a$. The radial gradient
there has norm tending to infinity as $r\downarrow a$.
\end{lemma}
\begin{proof}
The derivative is $\sqrt{B/(r-a)}$, and the second derivative is
$-\sqrt B/(2(r-a)^{3/2})$. The constants in
\eqref{eq:neck-adapter-constants} give the matching first jet. In a
radial interpolation, the second derivative may be smoothed through
negative functions while the first derivative remains positive; this
is the one-dimensional special case of Lemma~\ref{lem:smoothing}.
The radial Hessian eigenvalues $f''$ and $f'/r$ keep opposite signs.
On the square-root segment the radial gradient is strictly monotone
and becomes unbounded at $a$.
\end{proof}

\subsection{The two-sided extension theorem}
We now state explicitly what the global construction takes as input.
Neither the identity-return property nor a ruling is used in this
step. They were needed to create the protected band and its bendings,
not to extend a positive boundary jet.

\par\noindent\textbf{Formalization: pending.} See the precise object and dependency interfaces in the blueprint.\par
\begin{proposition}[Two-sided support completion]\label{prop:dual-radial-completion}
Suppose a compact saddle support annulus has injective source and
gradient maps, nested positively oriented source and gradient boundary
curves enclosing zero, and positive boundary traces as in
\eqref{eq:positive-traces}. Assume its outer boundary pair and its
reflected Legendre-dual inner pair obey the two visibility conditions
\eqref{eq:reflected-visible}. Then, after leaving an arbitrary
prescribed smaller protected annulus unchanged, it admits a saddle
support extension $\widetilde G$ on $\mathbb R^2\setminus\{0\}$ with
\begin{equation}\label{eq:radial-two-ends}
 \begin{aligned}
 \widetilde G(p)&=R_N|p|-\frac\mu2|p|^2+d_0
     && (0<|p|\ll1),\\
 \widetilde G(p)&=A|p|-\frac B{|p|}+d_\infty
     && (|p|\gg1),
 \end{aligned}
\end{equation}
where $R_N>A>0$ and $\mu,B>0$. The gradient is a
diffeomorphism
\begin{equation}\label{eq:full-gradient-annulus}
 \nabla\widetilde G:\mathbb R^2\setminus\{0\}
       \longrightarrow \{y\in\mathbb R^2:A<|y|<R_N\}.
\end{equation}
Every smoothing takes place outside the protected region.
\end{proposition}
\begin{proof}
\emph{The exterior end.}
Apply Lemma~\ref{lem:connector} to $(p_+,J\gamma_+)$, retaining an
incoming collar. Its terminal gradient is $Re_\theta$, and the
terminal source trace $P(\theta)$ is a positive Jordan curve with
$P\cdot e_\theta>0$. Apply the differential Legendre transform
$f=G^*$ on the gradient image. Its source seam is the circle
$y=Re_\theta$ and
\[
 f_r(Re_\theta)=P\cdot e_\theta>0,\qquad
 h''+Rf_r=R P_\theta\cdot Je_\theta>0.
\]
The last inequality is precisely positivity of the connector's
outgoing seam. Lemma~\ref{lem:quadratic-filling} therefore continues
$f$ radially inward, with injective gradient. Stop within the
quadratic radial region at $r=r_j$, and apply
Lemma~\ref{lem:neck-adapter}; choose its inner radius $a=A>0$.
The gradient of this dual continuation tends radially to infinity at
$r=A$. Inverse Legendre transformation consequently extends the
original potential toward $|p|=\infty$. Indeed, on the terminal
square-root germ $|p|=f_{\rm neck}'(r)$, so
$r=A+B/|p|^2$ and
\[
 p\cdot y-f_{\rm neck}(y)=A|p|-B/|p|-C.
\]
This is the second germ in \eqref{eq:radial-two-ends}.
For clarity, the three domains in this exterior construction are
$ p $-space for $G$, $y=\nabla G(p)$-space for $f=G^*$, and
$p=\nabla f(y)$-space when we return to the original potential.
Moving inward in the $y$-source annulus makes the $p$-gradient
circles expand. The square-root endpoint $|y|=A$ therefore becomes
the unbounded original-source end. Each relative smoothing retains
an open incoming collar and the exact terminal radial germ.

\emph{The interior end.}
Use reflected differential duality: put $K(x)=G^*(Sx)$. In terms
of the original inner traces $(p_-,\gamma_-)$, the new source and
gradient traces, with reversed parameter, are respectively
\[
 p^-(s)=S\gamma_-(-s),\qquad \gamma^-(s)=Sp_-(-s).
\]
The action period still vanishes, since
$\oint p_-\cdot d\gamma_-=-\oint\gamma_-\cdot dp_-=0$.
The original gradient curve $\gamma_-$ is the outer boundary
of the gradient annulus, so $S\gamma_-$ is the outer source boundary
for $K$. Reflection reverses planar orientation and $s\mapsto-s$
reverses it again; both new trace parametrizations are therefore
positive. The connector extends $K$ on its exterior side. After
taking $K^*$, this added region lies on the bounded source side of
the original $p_-$ boundary. This is the interior end we need. Its trace meets the same connector
hypotheses by \eqref{eq:reflected-visible}. Run the connector, then
take $f=K^*$ and use Lemma~\ref{lem:quadratic-filling} down to the
puncture. On the unchanged initial collar
\[
 K^*(y)=G(Sy),
\]
so the inverse reflection recovers the original $G$. The radial
quadratic germ of $K^*$ is invariant under $S$, giving the first
germ in \eqref{eq:radial-two-ends}. Its leading radius
$R_N=\mu R$ can be made arbitrarily large; we choose it larger than
$A$ and than the bounded intermediate gradient traces.

\emph{Global geometry and periods.}
We apply the degree criterion to different maps at different stages.
First, the connector source map $P$ has an embedded large outer
boundary enclosing its incoming source curve; its negative Jacobian
therefore proves source injectivity. Second, on the retained support
plus connector, the gradient has nested boundary images: the retained
original gradient boundary and the new radius-$R$ circle. Degree gives
a global inverse on this joined annulus, so the first dual potential
is single-valued.

Third, in each dual filling, retain the far incoming boundary in an
unchanged collar. The newly added small source circle is inside it,
whereas its gradient is a large circle enclosing the incoming gradient
trace. The coefficient in Lemma~\ref{lem:quadratic-filling} is chosen
after those traces are fixed. Degree makes the continued dual gradient
injective on every such compact annulus. The radial neck has strictly
monotone, expanding gradient radii, so the same argument persists down
to any circle above its inner endpoint. These are the global inverse
domains used to transform back. Relative smoothing changes no
retained boundary germ in these degree applications.

Finally, the original inner and outer extensions agree with $G$ on
open collars on their respective sides. They thus define one potential
$\widetilde G$ on the punctured plane. Its global gradient image is
verified separately by the following exhaustion. On the inner and outer radial collars,
for sufficiently small $\varepsilon>0$ and large $L$,
\[
 \nabla\widetilde G(\varepsilon e_\theta)
       =(R_N-\mu\varepsilon)e_\theta,
 \qquad
 \nabla\widetilde G(Le_\theta)
       =(A+B/L^2)e_\theta.
\]
The boundary circles are disjoint and nested, and the Hessian
Jacobian is negative on the intervening source annulus.
Lemma~\ref{lem:degree} proves that its gradient is a
diffeomorphism onto the intervening planar annulus. Exhausting the
punctured plane establishes \eqref{eq:full-gradient-annulus}.
All Legendre transformations are genuine differentials of globally
injective gradients on the relevant compact annuli; they therefore
introduce no missing periods. The inner and outer potentials recover
the original value and gradient on unchanged overlaps, so they
assemble into one globally defined support.
\end{proof}



\begin{figure}[p]
\centering
\includegraphics[width=\textwidth]{figures/05-completion.pdf}
\caption{Completion of the saved balanced numerical band and the proof's
domain changes. (a) The unchanged protected band is gold; the initial
outer and inner attachments are teal and purple.
(b) The blue patch begins the radial filling beyond the outer connector.
Both windows have the same physical bounds and scale.
(c) A positive exit feeds a visible connector with circular outgoing
gradient. Legendre duality turns that trace into a circular source seam
for quadratic radial filling. The two ends then receive the polar and
neck profiles described in Proposition~\ref{prop:dual-radial-completion}.
(d) A cutaway of the northern assembly, its reflection, and the gray
convex exterior, all under one common similarity. Gold location markers
are enlarged for visibility: the band is approximately
$1.17\times10^{-12}$ of the total height in this saved assembly.
The mesh matches first jets at internal interfaces and is shown before
relative smoothing; it does not certify a smooth tight torus.
Lemma~\ref{lem:smoothing} and Section~\ref{sec:model-saddle} supply the
smooth geometric completion in the proof.}
\label{fig:completion}
\end{figure}

\section{From the radial support to a tight torus}\label{sec:model-saddle}
Proposition~\ref{prop:dual-radial-completion} gives two explicit end
germs without passing through a logarithmic support. One end resolves
into a circular boundary, the other into a reflectable parabolic neck.
We then close the resulting negative-curvature annulus by a convex
annulus. This stage uses only the support geometry. It transports the
protected bending region by one fixed ambient isometry.

\subsection{Circular compactification and reflection}
Let $\widetilde G$ be the completed potential in
\eqref{eq:radial-two-ends}, and put $r=|p|$.
Using \eqref{eq:planar}, its image near the puncture is
\begin{equation}\label{eq:new-polar-immersion}
 X_{\widetilde G}(r,\theta)=
 ((R_N-\mu r)e_\theta,\ d_0+\tfrac\mu2r^2).
\end{equation}
Near infinity it is
\begin{equation}\label{eq:new-neck-immersion}
 X_{\widetilde G}(r,\theta)=
 ((A+B/r^2)e_\theta,\ d_\infty-2B/r).
\end{equation}
Apply the ambient isometry $T(x)=-x+d_\infty e_3$ to the whole
immersion, and shift the angular parameter by $\pi$. We use this same
normalization for the protected core and its bending fields: if their
original restrictions are $X|_V$ and $y|_V$, their restrictions after
normalization are $T\circ X|_V$ and $dT(y)|_V=-y|_V$. Thus subsequent
relative statements are literal in the normalized coordinates. At the
puncture, write $t=-r$;
then the northern boundary collar is
\begin{equation}\label{eq:ring-N}
 F_N(\theta,t)=(R_N+\mu t)e_\theta
       +\bigl(h-\tfrac\mu2t^2\bigr)e_3,
 \qquad -\delta<t\le0,
\end{equation}
where $h=d_\infty-d_0$. At the other end put $\eta=1/r$.
The neck is
\begin{equation}\label{eq:new-parabolic-neck}
 X(\theta,\eta)=(A+B\eta^2)e_\theta+2B\eta e_3,
 \qquad \eta\ge0,
\end{equation}
so its radius and height satisfy
\begin{equation}\label{eq:parabolic-profile}
 \rho_{\rm spatial}=A+\frac{z^2}{4B}.
\end{equation}
It is therefore smoothly reflected across $z=0$.

The puncture $p=0$ is \emph{not} filled by one surface point: it
resolves into a whole boundary circle. In \eqref{eq:ring-N}, the
angular and transverse derivatives at $t=0$ are
$R_NJe_\theta$ and $\mu e_\theta$, so the immersion is regular.
The neck \eqref{eq:new-parabolic-neck} is also regular at $\eta=0$,
where the corresponding derivatives are $AJe_\theta$ and $2Be_3$.

\par\noindent\textbf{Formalization: pending.} See the precise object and dependency interfaces in the blueprint.\par
\begin{proposition}[Completed saddle annulus]\label{prop:saddle}
The completed support compactifies, after reflection at its neck,
to a smooth embedded annulus $N$ with $K<0$ in its interior. Its
boundary circles have common radius $R_N$, at heights $\pm h$ with
$h>0$. Near them its collars are \eqref{eq:ring-N} and
\begin{equation}\label{eq:ring-S}
 F_S(\theta,t)=(R_N+\mu t)e_\theta
       +\bigl(-h+\tfrac\mu2t^2\bigr)e_3,
 \qquad -\delta<t\le0.
\end{equation}
On either collar,
\begin{equation}\label{eq:ring-curvature}
 K=\frac{t}{\mu(R_N+\mu t)(1+t^2)^2}.
\end{equation}
All bending fields supported in the protected region extend by zero
over the annulus, and vanish on full boundary collars.
\end{proposition}
\begin{proof}
\emph{Embeddedness of the upper half.}
By \eqref{eq:full-gradient-annulus} its horizontal projection is
globally injective on the punctured source. The two displayed end
germs show that projection extends to nested disjoint boundary
circles of radii $A<R_N$, with regular full surface collars.
Consequently the compactified upper half is embedded.

\emph{Height separation.}
The interior Gauss normal is represented by
$q(p)=(p,1)/\sqrt{1+|p|^2}$, up to orientation, and is never vertical
for $0<|p|<\infty$. Hence the vertical height has no interior
critical point. Formula \eqref{eq:new-parabolic-neck} gives strictly
positive height close to the lower neck circle, while
\eqref{eq:ring-N} gives height strictly below $h$ near the upper
circle. The maximum and minimum on the compact annulus therefore
occur at these boundary components. It follows that
\[
 h>0,\qquad 0<X_3<h\quad\text{in the upper interior}.
\]
The reflection at the even neck \eqref{eq:parabolic-profile} is
smooth; the two halves have disjoint interiors, so their union is
embedded and has negative curvature throughout the interior.

\emph{Boundary curvature and retained fields.}
For \eqref{eq:ring-N}, set $r(t)=R_N+\mu t$ and
$z(t)=h-\mu t^2/2$. A unit normal is
$(t e_\theta+e_3)/\sqrt{1+t^2}$; the metric coefficients are
$r(t)^2$ and $\mu^2(1+t^2)$, and the second fundamental
coefficients are $-r(t)t/\sqrt{1+t^2}$ and
$-\mu/\sqrt{1+t^2}$. This proves \eqref{eq:ring-curvature}.
Reflection gives the same intrinsic formula at the southern end.
Every change was outside the protected region, where all retained
bendings vanish on a collar; extension by zero is therefore smooth
and preserves their linearized metric equations.
\end{proof}

\subsection{Convex closure}\label{sec:model-convex}
The two boundary circles of the saddle annulus have radius $R_N$
and heights $\pm h$. We join them by a positive-curvature annulus
lying outside the cylinder of radius $R_N$. A concave radius as a
function of height gives a convex lateral surface. At each end we
prescribe the entire parabolic germ, so the final join is smooth.

\par\noindent\textbf{Formalization: proved.} See the precise object and dependency interfaces in the blueprint.\par
\begin{lemma}[Convex closure]\label{lem:convex}
For any $R_N,\mu,h>0$, there exists a convex body whose boundary
consists of the disks
\[
 D_\pm=\{(v,\pm h):|v|\le R_N\}
\]
and a smooth annulus $U_{\rm sp}$ with strictly positive curvature
in its interior. The annulus agrees with \eqref{eq:ring-N} and
\eqref{eq:ring-S} for $0\le t<\delta$, has outward Gauss map
\begin{equation}\label{eq:convex-gauss}
 \nu:U_{\rm sp}^{\circ}\longrightarrow
            \mathbb S^2\setminus\{\pm e_3\}
\end{equation}
a diffeomorphism, and satisfies
\begin{equation}\label{eq:convex-exterior}
 U_{\rm sp}^{\circ}\subset
      \{\rho_{\rm spatial}>R_N,\ |z|<h\},
 \qquad U_{\rm sp}\ne -U_{\rm sp}.
\end{equation}
\end{lemma}
\begin{proof}
Choose a strictly concave, even function $r_0$ on $[-h,h]$ with
$r_0(\pm h)=R_N$, $r_0>R_N$ in the interior, and near the endpoints
\begin{equation}\label{eq:meridian-radius}
 r_0(z)=R_N+\sqrt{2\mu(h-z)}\quad(z\text{ near }h),
 \qquad
 r_0(z)=R_N+\sqrt{2\mu(h+z)}\quad(z\text{ near }-h).
\end{equation}
For completeness, on $[0,h)$ choose a smooth function $v=-r_0'$
with $v'>0$, $v(z)=\epsilon z$ near zero, and
$v(z)=\sqrt{\mu/[2(h-z)]}$ near $h$. Choose the endpoint
neighborhoods small enough to interpolate with $v'>0$ between
these germs, and set
$r_0(z)=R_N+\int_z^h v(s)\,ds$ for $z<h$, with $r_0(h)=R_N$.
The endpoint singularity of $v$ is integrable. The expression is smooth
after even reflection at zero, and its lateral surface is smooth at
the endpoints in the $t$ coordinate used below. In the height
coordinate, $r_0$ is continuous on $[-h,h]$ and smooth on $(-h,h)$;
its endpoint slopes are infinite. It is strictly concave and agrees with
\eqref{eq:meridian-radius} at the upper endpoint.

Choose a nonzero smooth $b\ge0$ supported in $(0,h)$ and set
$r=r_0+\lambda b$. For sufficiently small $\lambda>0$ it remains
strictly concave, has unchanged end germs, and is no longer even.
The body $\mathcal K=\{(v,z):|z|\le h,\ |v|\le r(z)\}$ is convex,
since $r$ is concave. Its lateral annulus $(r(z)e_\theta,z)$ has
\[
 \nu=\frac{(e_\theta,-r'(z))}{\sqrt{1+r'(z)^2}},\qquad
 K=-\frac{r''(z)}{r(z)(1+r'(z)^2)^2}>0.
\]
As $z$ increases, $r'$ decreases from $+\infty$ to $-\infty$,
which gives the stated Gauss diffeomorphism. The upper equality in
\eqref{eq:meridian-radius}, together with
$z=h-\mu t^2/2$, gives $r=R_N+\mu t$ for $t\ge0$;
the southern equality works identically. Thus the joins preserve
full smooth germs. Since $r(z)>R_N$ for $|z|<h$, the lateral annulus
lies outside the prescribed cylinder, and non-evenness excludes
central symmetry. The disks close the auxiliary convex body. Smoothness
of its boundary across the disk rims is not needed: the smooth torus
is obtained by joining the lateral annulus to the saddle collars.
\end{proof}

The global gradient image in \eqref{eq:full-gradient-annulus} gives
$\rho_{\rm spatial}<R_N$ on $N^\circ$, and the height separation
in Proposition~\ref{prop:saddle} gives $|z|<h$ there. Hence
\[
 N\subset\mathcal C:=\{\rho_{\rm spatial}\le R_N,\ |z|\le h\},
\]
with strict inclusion for interior points. By
\eqref{eq:convex-exterior}, $U_{\rm sp}$ meets $N$ only along their
two boundary circles. Their parabolic germs coincide for negative
and positive $t$, producing a smooth embedded torus.
The curvature signs are negative on $N^\circ$, positive on
$U_{\rm sp}^\circ$, and zero exactly on the two joining circles.
By the Gauss area formula and \eqref{eq:convex-gauss},
\begin{equation}\label{eq:positive-integral}
 \int_{U_{\rm sp}^\circ}K\,dA=4\pi.
\end{equation}
Gauss--Bonnet gives total absolute curvature $8\pi$, hence
tightness. This uses the standard tightness characterization in
terms of total absolute curvature \cite{BK}.


\subsection{An asymmetric fixed region}
In this paper a \emph{marked} surface means that the unchanged
positive-curvature region has trivial ambient Euclidean stabilizer.
This geometric condition rules out congruence of images; it does not
merely restrict allowed parametrizations. For a set $A\subset\R^3$,
$\Stab_{\Isom(\R^3)}A$ is the group of Euclidean isometries carrying
$A$ onto itself.

\par\noindent\textbf{Formalization: pending.} See the precise object and dependency interfaces in the blueprint.\par
\begin{lemma}[Small affine marking]\label{lem:ellipses-new}
For every sufficiently small $\epsilon>0$, the map
\[
 A_\epsilon=\begin{pmatrix}
 1&0&\epsilon\\0&1+\epsilon&2\epsilon\\0&0&1
 \end{pmatrix}
\]
tends to $\Id$ and satisfies
$\Stab_{\Isom(\R^3)}A_\epsilon(U_{\rm sp}^\circ)=\{\Id\}$.
\end{lemma}
\begin{proof}
An isometry preserving the open annulus preserves its closure and permutes
its two boundary ellipses. Their centers are
$\pm h(\epsilon,2\epsilon,1)$; their planes are horizontal; their unequal
semiaxes lie along $e_1,e_2$ and have lengths $R_N,(1+\epsilon)R_N$.
The midpoint of the centers is fixed, so the translation is zero.
The planes and principal axes force the orthogonal part to be diagonal
with entries $\pm1$. If both centers are fixed, all entries are $+1$,
because each coordinate of $(\epsilon,2\epsilon,1)$ is nonzero.
If they are exchanged, all are $-1$. This would give central inversion
of the annulus and hence $U_{\rm sp}=-U_{\rm sp}$, contrary to
Lemma~\ref{lem:convex}. Only the identity remains.
\end{proof}

\subsection{The completed protected module}\label{sec:realization}
\par\noindent\textbf{Formalization: pending.} See the precise object and dependency interfaces in the blueprint.\par
\begin{theorem}[Protected identity-return realization]\label{thm:protected-realization}
There exist a smooth torus $M$, a tight embedding $X:M\to\R^3$, a
compact annulus $V\Subset\{K_X<0\}$, an open annular neighborhood
$W\supset V$, and a nonempty interval $I$ such that:
\begin{enumerate}[label=\textup{(\roman*)}]
\item The Gauss map and horizontal projection are injective on $W$.
In planar support coordinates, $\det D^2G<0$ and both the source and
gradient maps are diffeomorphisms onto planar annuli.
\item Under the fixed source identification with the protected ruled
band, the transported field $\bar y_F=-y_F$, with $y_F$ given by
\eqref{eq:yF}, defines an injective linear map from $C_c^\infty(I)$
to bendings supported in $V$. It is extended by zero outside that band.
Disjoint profile supports give disjoint bending supports.
\item Once the balanced speed is fixed and the core has undergone the
fixed ambient normalization of Section~\ref{sec:model-saddle}, the
completion leaves a prescribed compact subannulus unchanged. Every sufficiently small $L^1$
speed change preserving closure, period balance, and visibility admits
the same type of completion, with its own protected core.
\item The embedding has no open planar subset and total absolute curvature
$8\pi$.
\end{enumerate}
\end{theorem}
\begin{proof}
\emph{Build and protect the local module.}
Lemma~\ref{lem:seed} and Theorem~\ref{thm:spike} give a balanced embedded
frame with visible traces. Equation \eqref{eq:frame-speed} gives $I_0=0$,
so Theorem~\ref{thm:ruled} provides the kernel. Shrink the band for Gauss
and projection injectivity, and choose the protected level annulus $V_0$.
Enlarge it slightly to a compact $V$, still separated from both exit
collars. All fields from $C_c^\infty(I)$ are supported there.

\emph{Complete outside the protected neighborhood.}
The central mixed-entry formula and exit lemma give two positive seams.
The connector and quadratic radial filling lemma, combined with
Legendre duality and the neck adapter, give the two-sided support
of Proposition~\ref{prop:dual-radial-completion}.
Proposition~\ref{prop:saddle} yields the embedded saddle annulus, and
Lemma~\ref{lem:convex} closes it into an embedded torus. Apart from the
one ambient normalization applied uniformly to the core and fields,
every change is outside a neighborhood of $V$, so the transported
bendings extend smoothly by zero.
The positive curvature integral is $4\pi$; Gauss--Bonnet gives $8\pi$
absolute curvature and tightness. Curvature vanishes only on two circles,
so no planar open set exists.

\emph{Record the relative scope of the construction.}
Small $L^1$ speed changes give small positional changes with tangent
directions fixed. Embedding and visibility persist as in
Theorem~\ref{thm:spike}. For each final balanced speed we choose its thin
band and exit collars and repeat the completion, leaving its prescribed
normalized core fixed. This asserts existence of the same type of completion, not
an identification of different cores or a common exterior for all speeds.
\end{proof}

\par\noindent\textbf{Formalization: pending.} See the precise object and dependency interfaces in the blueprint.\par
\begin{corollary}[Marked protected module]\label{cor:marked-module}
There are a tight torus embedding $X:M\to\R^3$, a compact
$V\Subset\{K_X<0\}$, an interval $I\ne\varnothing$, and an injective
linear map $\mathcal Y:C_c^\infty(I)\to C^\infty(M,\R^3)$ whose values
are bendings supported in $V$, preserving disjointness of supports, such
that
\[
 \Stab_{\Isom(\R^3)}X(\{K_X>0\})=\{\Id\}.
\]
The torus has no open planar subset and total absolute curvature $8\pi$.
\end{corollary}
\begin{proof}
Apply $A_\epsilon$ to the preceding torus and transport fields by
\eqref{eq:affine-transport}. On the protected band the resulting map is
$\mathcal Y(F)=A_\epsilon^{-T}\bar y_F=-A_\epsilon^{-T}y_F$,
using the same source identification, and is extended by zero.
This preserves embedding, curvature signs,
tightness, strain, and supports. Lemma~\ref{lem:ellipses-new} marks the
positive-curvature region. The equality case of \eqref{eq:tight-bound}
gives the same absolute curvature value after the affine change.
\end{proof}

\section{Exact metric branching}\label{sec:reconstruction}
Fix the marked tight torus from Corollary~\ref{cor:marked-module}.
Its reference metric is $g_X$. The sign branches constructed in this
section induce a new common metric $g$, given by a positive quadratic
correction to $g_X$. We first make one exact pair, then choose finitely
or countably many independent signs. Tightness follows from the
unchanged positive-curvature region, and image noncongruence follows
from its marking together with intrinsic isometry rigidity.

\subsection{Finite independent signs}
The metric reconstruction map has the exact quadratic expansion
\begin{equation}\label{eq:metricquad}
 \Gmap(X+y)=\Gmap(X)+S_Xy+\langle dy,dy\rangle.
\end{equation}
\par\noindent\textbf{Formalization: proved.} See the precise object and dependency interfaces in the blueprint.\par
\begin{lemma}[Disjoint quadratic branching]\label{thm:branch}
If $y_1,\ldots,y_m\in\ker S_X$ have pairwise disjoint supports, then
\begin{equation}\label{eq:abstractsign}
 g_{X+\sum_{j=1}^m\varepsilon_j a_jy_j}
 =g_X+\sum_{j=1}^m a_j^2\langle dy_j,dy_j\rangle,
 \qquad\varepsilon_j\in\{-1,1\}.
\end{equation}
The same identity holds for an infinite series converging in $C^\infty$.
\end{lemma}
\begin{proof}
Every linear term vanishes by the strain equation. A derivative is zero
off its field's support, so disjoint supports kill every mixed quadratic
term pointwise. This proves the finite formula. Smooth convergence allows
passage to the limit in the metric and its derivatives.
\end{proof}

Two separate issues remain: tightness of the perturbed embeddings, and
noncongruence of their images. Both are controlled by the unchanged region.
\par\noindent\textbf{Formalization: partial.} See the precise object and dependency interfaces in the blueprint.\par
\begin{proposition}[Relative preservation of absolute curvature]\label{prop:budget}
Let $X$ be a closed embedded surface and $V\Subset\{K_X<0\}$.
Every sufficiently small $C^2$ perturbation supported in $V$ remains an
embedding and has the same total absolute curvature. If $X$ is tight,
so is the perturbation.
\end{proposition}
\begin{proof}
Compactness and smallness preserve negative curvature on $V$ and preserve
embedding. Off $V$ the map is unchanged, so the positive-curvature region
$P$ and its curvature integral are unchanged. For either surface,
Gauss--Bonnet gives
\[
 \int_M|K|\,dA=2\int_P K\,dA-2\pi\chi(M).
\]
Thus the absolute curvature is unchanged, and equality in
\eqref{eq:tight-bound} preserves tightness.
\end{proof}

\par\noindent\textbf{Formalization: partial.} See the precise object and dependency interfaces in the blueprint.\par
\begin{lemma}[An intrinsic isometry fixed on an open set]\label{lem:fixed-open}
An isometry of a connected smooth Riemannian manifold which fixes a
nonempty open set pointwise is the identity. The same holds for a connected
manifold with boundary when the fixed set meets its interior.
\end{lemma}
\begin{proof}
In the interior, the set of points fixed with identity differential is
nonempty and closed. It is open by the exponential-map relation
$f(\exp_pv)=\exp_{f(p)}(df_pv)$. Connectedness propagates the identity.
For a manifold with boundary, apply this to its connected interior and
then use continuity at the boundary.
\end{proof}

\par\noindent\textbf{Formalization: pending.} See the precise object and dependency interfaces in the blueprint.\par
\begin{proposition}[Finite sign families]\label{prop:finite-sign}
Let $X:M\to\R^3$ be a tight embedding of a closed connected surface,
with $P=\{K_X>0\}$ and $\Stab_{\Isom(\R^3)}X(P)=\{\Id\}$.
Let $y_1,\ldots,y_m$ be nonzero bendings with disjoint supports in one
compact $V\Subset\{K_X<0\}$. For all sufficiently small positive
$a_j$, the $2^m$ maps in \eqref{eq:abstractsign} are tight embeddings
with a common metric, agree off $V$, and have pairwise noncongruent images.
\end{proposition}
\begin{proof}
The quadratic lemma gives the common metric and the curvature proposition
gives tight embeddings. Suppose an ambient isometry carries one image
onto another. Gaussian curvature is preserved, so it must carry the
unchanged positive region $X(P)$ onto itself. The stabilizer hypothesis
makes that ambient isometry the identity.

The images therefore coincide. If their parametrizations are $X_\varepsilon$
and $X_\delta$, then $\varphi=X_\delta^{-1}\circ X_\varepsilon$ is an
isometry of the common intrinsic metric. It fixes $P$ pointwise because
both maps equal $X$ there. The fixed-open-set lemma gives $\varphi=\Id$.
Thus the maps coincide pointwise. On a point where $y_j\ne0$, all other
fields vanish, so $\varepsilon_j=\delta_j$. This holds for every $j$.
\end{proof}
For $m=1$, this is already the exact nonuniqueness assertion. Increasing
$m$ repeats the same localized mechanism independently.


Figure~\ref{fig:cantor} sketches the passage from independent local
supports to the infinite sign space. Disjoint intervals in the leaf
coordinate give disjoint bending supports on the surface. The metric
sees the square of each amplitude and loses its sign. Rapid decay in
every derivative order makes the infinite series smooth, including at
an accumulation of the supports. The fixed marked region then separates
the resulting congruence classes.
\begin{figure}[htbp]
\centering
\includegraphics[width=\textwidth]{figures/06-cantor-family.pdf}
\caption{The Cantor family as independent sign choices.
(a) Schematic weighted test profiles on disjoint transverse intervals
$J_j\Subset I$ accumulating at an interior leaf label. Their supports
select disjoint annuli for $y_j=\mathcal YF_j$.
The plotted heights only illustrate decay; the proof chooses amplitudes
controlling all required $C^k$ norms.
(b) Three finite levels show $2$, $4$, and $8$ choices; continuing gives
$\{-1,1\}^{\mathbb N}$ with the product topology.
Disjoint supports eliminate mixed metric terms, and therefore every
branch has the same metric
$g_X+\sum_j a_j^2\langle dy_j,dy_j\rangle$.
Theorem~\ref{thm:localized-sign} proves smooth convergence, tightness,
noncongruence, and the embedding into the metric fiber modulo Euclidean
isometries. This totally disconnected parameter family is not a
continuous isometric flex.}
\label{fig:cantor}
\end{figure}


\subsection{Countably many signs and the quotient topology}
\par\noindent\textbf{Formalization: pending.} Deferred from the first-pair objective.\par
\begin{lemma}[Hausdorff Euclidean orbit space]\label{lem:orbit-space}
For a closed Riemannian manifold $(M,g)$,
$\Emb_g(M,\R^3)/\Isom(\R^3)$ is metrizable and hence Hausdorff.
\end{lemma}
\begin{proof}
Subtract the mean of each embedding with respect to the Riemannian
volume $d\mu_g$. Centering removes
translations, so the quotient is the quotient of centered embeddings by
$O(3)$. On that space use the invariant metric
\[
 d(f,h)=\sum_{k\ge0}2^{-k-1}\min\{1,\|f-h\|_{C^k(g)}\}.
\]
Since $O(3)$ is compact, the orbit distance
$\bar d([f],[h])=\min_{Q\in O(3)}d(f,Qh)$ is a metric. Invariant metric
balls show that it induces the quotient topology.
\end{proof}

\par\noindent\textbf{Formalization: partial.} Deferred from the first-pair objective.\par
\begin{theorem}[Localized sign families on a marked tight surface]
\label{thm:localized-sign}
Let $X:M\to\R^3$ be a tight embedding of a closed connected surface
whose positive-curvature image has the trivial ambient stabilizer
required in Proposition~\ref{prop:finite-sign}. Let $y_j$, $j\ge1$,
be nonzero bendings with pairwise disjoint supports in
one compact $V\Subset\{K_X<0\}$. Choose positive amplitudes such that
\[
 X_\varepsilon=X+\sum_{j\ge1}\varepsilon_ja_jy_j,
 \qquad\varepsilon\in\{-1,1\}^{\mathbb N},
\]
converges smoothly uniformly in the signs and is sufficiently small in
$C^2$. Then the maps are tight embeddings with the common smooth metric
\[
 g=g_X+\sum_{j\ge1}a_j^2\langle dy_j,dy_j\rangle.
\]
Their images are pairwise noncongruent, and the sign map induces a
topological embedding into $\Emb_g(M,\R^3)/\Isom(\R^3)$.
Such amplitudes always exist.
\end{theorem}
\begin{proof}
\emph{Choose a uniformly smooth and small series.}
For a prescribed small $\epsilon>0$, require
\[
 a_j\max_{0\le k\le j+2}\|y_j\|_{C^k}<\epsilon2^{-j}.
\]
For each derivative order the tail is uniformly summable, independently
of the signs. The full sum is also uniformly small in $C^2$.
The quadratic identity and the curvature proposition give the common
metric and tight embeddings, with the positive region unchanged.

\emph{Rule out congruence and reparametrization.}
The argument in Proposition~\ref{prop:finite-sign} applies verbatim:
an ambient congruence preserves $X(P)$ and is the identity; the resulting
intrinsic reparametrization fixes $P$ and is the identity. At a point
where $y_j\ne0$, disjointness then determines its sign. Thus every sign
sequence gives a different congruence class.

\emph{Verify the Cantor topology.}
Agreement of sufficiently many initial signs makes the difference of two
maps a uniformly small tail in each $C^k$ norm. Hence the sign map is
continuous for the product topology. Its domain is compact and its
quotient target is Hausdorff by Lemma~\ref{lem:orbit-space}. A continuous
injection from a compact space into a Hausdorff space is a topological
embedding.
\end{proof}

\section{Consequences of the protected band}\label{sec:applications}
\subsection{The tight metric fiber}\label{sec:tight-main}
\begin{proof}[Proof of Theorem~\ref{thm:main-fiber}]
Use the marked module in Corollary~\ref{cor:marked-module}. Choose disjoint
intervals $J_j\Subset I$ and nonzero $F_j\in C_c^\infty(J_j)$, and let
$y_j=\mathcal YF_j$. The module preserves disjoint supports. Choose the
amplitudes as in Theorem~\ref{thm:localized-sign}. It gives a common metric,
tight embeddings, noncongruent images, and the Cantor embedding into the
quotient. All maps equal $X$ off $V$, in particular on the same nonempty
positive-curvature open set.

The reference curvature vanishes only on the two joining circles.
The perturbed curvature remains negative on $V$ and is unchanged elsewhere,
so no member has an open planar subset. Tightness gives total absolute
curvature $8\pi$. Selecting any two distinct sign sequences proves the
first assertion with the same metric. The finite proposition also proves
the two-realization phenomenon directly from one nonzero field.
\end{proof}

\par\noindent\textbf{Formalization: pending.} Deferred from the first-pair objective.\par
\begin{corollary}[Smooth tight nonrigidity]\label{cor:tight-nonrigidity}
Smooth embedded tight tori are not determined up to congruence by their
intrinsic metrics. One such metric has a Cantor family of noncongruent
realizations.
\end{corollary}
\begin{proof}
Apply Theorem~\ref{thm:main-fiber}.
\end{proof}

\subsection{Closed asymptotic curves and zero linking}
\begin{proof}[Proof of Theorem~\ref{thm:main-annulus}]
Take $N$ from Proposition~\ref{prop:saddle}. The collar formulas give
its circular tangent-plane boundaries and negative interior curvature.
The northern Gauss coordinates cover the northern hemisphere with its
pole removed. Reflection covers the southern hemisphere with its pole
removed; the smooth neck maps to the equator. Each half is parametrized
by its normal, and the neck has nonzero curvature. Together these give
a Gauss diffeomorphism onto the twice-punctured sphere.

The protected leaves in \eqref{eq:closed-leaves} remain embedded and
asymptotic. To compute their normal linking, cap $N$ by the planar disks
$D_\pm$ at heights $\pm h$. The height bounds show that
$S=N\cup D_+\cup D_-$ is an embedded topological sphere.
A small normal push-off of a protected leaf is disjoint from $S$:
near the leaf this follows from a tubular neighborhood, and away from
the leaf from compact separation. The leaf bounds a disk in $S$ disjoint
from the push-off. Intersection with that disk computes linking zero.
\end{proof}

\par\noindent\textbf{Formalization: pending.} Deferred from the first-pair objective.\par
\begin{corollary}[Degenerate asymptotic return on a tight torus]
\label{cor:HK-degenerate}
A smooth embedded tight torus contains an open annulus of closed
asymptotic curves in its negative-curvature region such that every leaf
satisfies
\[
 \int_\Gamma k_gk_n|K|^{-1/2}\,ds=0.
\]
Here $k_g=\langle\nabla_TT,Z\rangle$ and $k_n=\II(Z,Z)$, with $Z$
a unit tangent perpendicular to the oriented leaf tangent $T$.
\end{corollary}
\begin{proof}
Use the unmarked torus in Theorem~\ref{thm:protected-realization}.
Its protected return is the identity, so each leaf has multiplier one.
Lemma~\ref{lem:holonomy-integral-new} gives the integral's vanishing.
\end{proof}

\par\noindent\textbf{Formalization: pending.} Deferred from the first-pair objective.\par
\begin{corollary}[Fixed-boundary isometric saddle annuli]
\label{cor:fixed-boundary-annuli}
There are two noncongruent smooth embedded annuli with strictly negative
Gaussian curvature, the same induced metric, and the same two convex
planar boundary curves. They agree on full boundary collars.
\end{corollary}
\begin{proof}
\emph{Choose strictly negative boundaries and eliminate their symmetries.}
Cut the northern and southern collars of $N$ at $t=-\eta_+$ and
$t=-\eta_-$, where $0<\eta_+\ne\eta_-\ll1$. Retain the protected
support. These cuts give strictly negative curvature even at the boundary
and distinct circular radii. Apply $A_\epsilon$ from
Lemma~\ref{lem:ellipses-new} and transport the bendings by
\eqref{eq:affine-transport}. The new boundary curves are noncircular
ellipses in parallel planes with different semiaxes, so an isometry
preserving their union must preserve each individually.

Their principal axes force the orthogonal part of such an isometry to
be diagonal. Their centers are distinct multiples of
$(\epsilon,2\epsilon,1)$. Fixing both centers forces all three diagonal
signs to be $+1$, and then the translation is zero. Thus their union
has trivial ambient stabilizer.

\emph{Construct the exact pair and exclude identical images.}
Let $\widetilde X$ be the truncated affine embedding and $y\ne0$ a
retained bending. For small $t>0$, $\widetilde X\pm ty$ are embedded
strictly negative-curvature annuli. They agree on full boundary collars,
and \eqref{eq:metricquad} gives a common metric.
Any ambient congruence must preserve the boundary union and is therefore
the identity. If the images coincided, their induced reparametrization
would be an intrinsic isometry fixing the common collar pointwise.
Lemma~\ref{lem:fixed-open} would make it the identity on the whole annulus,
forcing $y=0$. This contradiction proves noncongruence.
\end{proof}

\subsection{Complete negatively curved ends}\label{sec:complete-cylinder}
For a complete cylinder we cut the saddle annulus before its zero-curvature
boundary circles and attach strictly convex radial profiles. Agreement of
full germs at the cuts preserves smoothness and the compact bending module.

\par\noindent\textbf{Formalization: proved.} Deferred from the first-pair objective.\par
\begin{lemma}[Strictly convex continuation of a radial profile]
\label{lem:complete-profile}
Suppose $q_0$ is smooth near $H$, with $q_0(H)>0$, $q_0'(H)>0$,
and $q_0''>0$ there. For every $\beta>0$ there is a smooth $q$ on
$[H,\infty)$ equal to $q_0$ near $H$, with $q,q',q''>0$, and with
$q''=\beta$ eventually. In particular $q'(z)\to\infty$.
\end{lemma}
\begin{proof}
Choose $d>0$ with $q_0''>0$ on $[H,H+2d]$. A smooth cutoff from zero
near $H$ to one for $z\ge H+d$ gives a positive smooth function $b$
equal to $q_0''$ near $H$ and to $\beta$ for $z\ge H+d$.
Define
\begin{equation}\label{eq:complete-profile}
 q(z)=q_0(H)+q_0'(H)(z-H)+\int_H^z(z-v)b(v)\,dv.
\end{equation}
Then $q''=b>0$, so $q'$ and $q$ stay positive. Near $H$ the second
derivative and initial first jet agree with those of $q_0$, giving full
agreement there. Eventually $q$ is quadratic with positive second derivative.
\end{proof}

\begin{proof}[Proof of Theorem~\ref{thm:complete-cylinder}]
\emph{Step 1: cut at strictly negatively curved circles.}
Use the unmarked $N$ of Proposition~\ref{prop:saddle} and choose sufficiently small $0<\eta<\delta$ so that cutting both
collars at $t=-\eta$ retains the protected support and $H>0$. Put
\begin{equation}\label{eq:complete-cut}
 R_\eta=R_N-\mu\eta>0,\qquad
 H=h-\frac\mu2\eta^2>0.
\end{equation}
The retained annulus $N_\eta$ has boundary circles at heights $\pm H$
and radius $R_\eta$, with $K<0$ everywhere up to those boundaries.
Its interior lies between these heights because the height function has
no interior critical point.

\emph{Step 2: extend the circular germs to convex ends.}
The northern collar, solved for radius as a function of height, is
\begin{equation}\label{eq:complete-collar}
 \begin{gathered}
 q_0(z)=R_N-\sqrt{2\mu(h-z)},\\
 q_0'(z)=\frac{\mu}{\sqrt{2\mu(h-z)}},\qquad
 q_0''(z)=\frac{\mu^2}{[2\mu(h-z)]^{3/2}}.
 \end{gathered}
\end{equation}
At $z=H$ one has $q_0=R_\eta$, $q_0'=1/\eta>0$, and
$q_0''=1/(\mu\eta^3)>0$.
The continuation lemma gives $q$ with the same germ and eventually
quadratic growth. Attach
\begin{equation}\label{eq:complete-ends}
 E_\pm(\theta,z)=q(z)e_\theta\pm ze_3,
 \qquad z\in[H,\infty).
\end{equation}
Full-germ agreement makes both joins smooth. The resulting source is
$S^1\times\R$ with no boundary.

\emph{Step 3: verify curvature, embedding, and completeness.}
For the upper end,
\begin{equation}\label{eq:complete-end-metric}
 g=q^2d\theta^2+(1+(q')^2)dz^2,\qquad
 n_+=\frac{-e_\theta+q'e_3}{\sqrt{1+(q')^2}}.
\end{equation}
Its second fundamental coefficients are $q/\sqrt{1+(q')^2}$ and
$-q''/\sqrt{1+(q')^2}$, with zero mixed coefficient, so
\begin{equation}\label{eq:complete-end-curvature}
 K=-\frac{q''}{q(1+(q')^2)^2}<0.
\end{equation}
Reflection gives the same result below, including at both smooth seams.
Height determines $z$ on each end and its positive radius determines
$\theta$, so each end is embedded. Their interiors lie strictly above
or below the core's height range, excluding all intersections.

The embedding is proper: a compact ambient set bounds height and hence
bounds the end parameters, leaving only a compact part of the source.
An intrinsic Cauchy sequence has bounded Euclidean image, since induced
distance dominates ambient distance. Properness gives a subsequence
converging on the surface, and local smooth coordinates turn this into
intrinsic convergence. The Cauchy sequence therefore converges, proving
completeness.

\emph{Step 4: identify the Gauss map and curvature integral.}
At the upper join $n_+=(-\eta e_\theta+e_3)/\sqrt{1+\eta^2}$,
matching the core normal. Along the end,
\[
 \frac d{dz}\frac{q'}{\sqrt{1+(q')^2}}
 =\frac{q''}{(1+(q')^2)^{3/2}}>0,
 \qquad\frac{q'}{\sqrt{1+(q')^2}}\longrightarrow1.
\]
Writing $\zeta_\eta=(1+\eta^2)^{-1/2}$, the upper end covers the
punctured cap $\zeta_\eta\le n_3<1$ once, and reflection covers
$-1<n_3\le-\zeta_\eta$ below. The retained core covers
$|n_3|\le\zeta_\eta$ once by
Theorem~\ref{thm:main-annulus}. The descriptions agree at the seams,
where $K\ne0$, so the global normal is a diffeomorphism onto the
twice-punctured sphere. Its area gives $\int_\Sigma|K|\,dA=4\pi$.
For comparison, each end contributes
\[
 2\pi\int_H^\infty\frac{q''}{(1+(q')^2)^{3/2}}\,dz
 =2\pi(1-\zeta_\eta).
\]

\emph{Step 5: retain the characteristic and bending module.}
The surface is unchanged near the protected band, so its closed leaves,
identity return, and normal fields along the leaves are unchanged.
Their linking remains zero. All bendings vanish near the cuts and extend
by zero to the ends, retaining their strain equations, injectivity, and
common compact support.
\end{proof}
The eventual quadratic profiles have $K\to0$ at infinity. Completeness
therefore comes with strict negativity everywhere, but no uniform negative
upper bound. The result addresses the embedded-existence question in
\cite[after Theorem~3.4]{NiZhangZhang}; their plane-topology obstruction
does not apply to a cylinder.

\appendix
\section{General ruled bands and projective return}\label{sec:general-ruled}
This appendix generalizes the principal-normal theory. It is not needed
for the explicit core or its completion. The general return has two
projective parameters rather than the single period $I_0$, and the bending
equation transports a transverse vector field rather than a scalar alone.

\par\noindent\textbf{Formalization: proved.} Deferred from the first-pair objective.\par
\begin{proposition}[Riccati return and a two-jet criterion]
\label{prop:general-ruled-return}
Let $X(t,u)=c(t)+ue(t)$ with $|e|=1$, smooth $L$-periodic $c,e$, and
$\Delta=\det(c',e,e')\ne0$. Suppose $u=0$ is a closed nonruling
asymptotic leaf. Then the nonruling characteristics satisfy
\begin{equation}\label{eq:general-riccati}
 u'=q(t,u)=q_1(t)u+q_2(t)u^2,
 \qquad q=-\frac{\langle c''+ue'',(c'+ue')\times e\rangle}{2\Delta}.
\end{equation}
With
\[
 A(t)=\int_0^tq_1(s)\,ds,\quad \mu=e^{A(L)},\quad
 \nu=\int_0^Lq_2(t)e^{A(t)}\,dt,
\]
the local return is
\begin{equation}\label{eq:general-return}
 \mathcal R(u)=\frac{\mu u}{1-\nu u}.
\end{equation}
Identity return is equivalent to $\mu=1$, $\nu=0$, or equivalently
$\mathcal R'(0)=1$ and $\mathcal R''(0)=0$.
\end{proposition}
\begin{proof}
Put $W=|(c'+ue')\times e|$. The mixed second fundamental coefficient
is $\Delta/W$ and the $uu$ coefficient is zero. The $tt$ numerator
is the quadratic polynomial in \eqref{eq:general-riccati}, whose constant
term vanishes because $u=0$ is asymptotic. This gives the characteristic
equation and $K=-\Delta^2/W^4<0$. Substitution $u=e^{A(t)}z$ gives
$z'=q_2e^Az^2$, and integration yields the displayed return.
\end{proof}
Without a chosen closed leaf there is also a constant term $q_0(t)$.
The Riccati equation is the projectivization $u=x/y$ of
\[
 \binom{x}{y}'=
 \begin{pmatrix}q_1/2&q_0\\-q_2&-q_1/2\end{pmatrix}\binom{x}{y}.
\]
Its monodromy, the linear solution map over one period, is in
$\mathrm{SL}(2,\R)$; projective identity means that
matrix is $\pm\Id$. To obtain an affine band one must also have a finite
periodic leaf, since a projective solution can pass through $u=\infty$.
Near the finite leaf in the proposition this issue does not arise.
For principal-normal rulings, $q_1=-\lambda/2$, $q_2=-D/2$, so
$\mu=1$ automatically and $\nu=-\rho(0)I_0$.

\par\noindent\textbf{Formalization: pending.} Deferred from the first-pair objective.\par
\begin{theorem}[Ruled localization as transverse transport]
\label{thm:general-ruled-kernel}
Under the preceding hypotheses, a sufficiently thin one-sided band
$0<u<b$ admits a nonzero compactly supported bending exactly when the
return is the identity. In the identity case, the supported kernel on
any saturated subband is parametrized by one arbitrary compactly supported
smooth function of its leaf parameter.
\end{theorem}
\begin{proof}
\emph{Reconstruct the field from one scalar.}
The $uu$ equation and support imply $\langle e,y\rangle=0$.
For $p=\langle X_t,y\rangle$, the mixed equation becomes
$\langle e',y\rangle=p_u/2$. Thus
\begin{equation}\label{eq:general-bending-reconstruction}
 \langle e,y\rangle=0,\qquad \langle X_t,y\rangle=p,
 \qquad\langle e',y\rangle=p_u/2
\end{equation}
uniquely determines $y$, since $\Delta\ne0$. The remaining equation is
\[
 p_t+qp_u=\left(q_u+\tfrac12(\log|\Delta|)'\right)p.
\]
For an explicit verification, write $e'=vf$, $n=e\times f$,
$c'=\alpha e+\beta f+\gamma n$, and $f'=-ve+\omega n$.
Then $\Delta=v\gamma$, and the right-hand coefficient is
$(\beta\omega+\gamma'+uv\omega)/\gamma$, which agrees with the formula.
Setting $h=p/\sqrt{|\Delta|}$ gives
\begin{equation}\label{eq:transverse-vector-transport}
 h_t+qh_u=q_uh,
 \qquad[\partial_t+q\partial_u,h\partial_u]=0.
\end{equation}

\emph{Identify the return condition.}
If $u=\phi_t(v)$ is the characteristic flow, then
\[
 h(t,\phi_t(v))=\partial_v\phi_t(v)h(0,v),\qquad
 h(0,\mathcal R(v))=\mathcal R'(v)h(0,v).
\]
Thus the transported object is the transverse vector field $h\partial_u$.
For identity return any compactly supported initial function on a
complete-leaf interval descends smoothly. Reconstruction gives every
supported bending on that saturated subband.

If return is not the identity, shrink $b$ to remove all fixed points in
$(0,b)$: the displayed projective return has at most one other fixed
point besides zero. A characteristic carrying nonzero $h$ remains nonzero
and must stay in a compact subset of the band. Its successive returns
are a monotone sequence in a compact interval, and their limit would be
an interior fixed point. This contradiction forces $h=0$ and hence $y=0$.
\end{proof}
Three distinct nearby closed nonruling leaves crossing a common ruling
once force a general projective return to be the identity. In the
principal-normal class the central leaf and one additional leaf suffice.
These finite tests depend on the projective form and do not apply to
arbitrary smooth surface foliations.

\section{Verification of the corrugated seed}\label{app:seed-estimates}
We prove the package in Lemma~\ref{lem:seed}. The estimates serve only
to give uniform visibility margins. The root and action calculations
provide closure, while visibility will also imply embedding of the
partner curve.

\subsection{Visibility implies embedding of a positive tangent partner}
\par\noindent\textbf{Formalization: pending.} See the precise object and dependency interfaces in the blueprint.\par
\begin{lemma}[Visible positive tangent partners]\label{lem:visible-partner}
Let $\beta$ be a regular positively oriented Jordan curve, and let
$\delta$ be closed with $\delta'=m\beta'$, $m>0$. If $|\delta|>R$
and $\beta'\cdot V_R(\delta)>0$, then $\delta$ is a positively oriented
Jordan curve enclosing the disk $|z|\le R$.
\end{lemma}
\begin{proof}
The unit tangent $U=\beta'/|\beta'|$ has degree one by the turning-number
theorem. Its positive scalar product with $V_R(\delta)$ makes those
circle-valued maps homotopic, so $V_R(\delta)$ has degree one.
Write $\delta=\rho e^{i\alpha}$ on the universal cover. An angle for the
visibility vector is $\varphi=\alpha+\arccos(R/\rho)$, whose second
summand is periodic. Thus $\delta$ winds once around zero. Moreover
\begin{equation}\label{eq:visibility-embedding}
 \varphi'=\frac{\delta'\cdot V_R(\delta)}{\sqrt{\rho^2-R^2}}>0.
\end{equation}
The angle increases strictly by $2\pi$ in one period. Equal positions at
distinct parameters would give equal angles modulo $2\pi$, a contradiction.
Hence $\delta$ is embedded and positively oriented. The disk is connected,
contains zero, and misses the curve, so it lies in the bounded component.
\end{proof}

\subsection{Root, first curve, covariance, and action}
For fixed $N$, consider the probability density proportional to
$(1+\epsilon\sin x)^2e^{-k\cos x}$. Equation
\eqref{eq:corrugated-root} is $\mathbb E_k[\cos x]=-1/2$, and
\[
 \frac d{dk}\mathbb E_k[\cos x]=-\operatorname{Var}_k(\cos x)<0.
\]
Concentration at the extrema of $\cos x$ gives limits $1$ as
$k\to-\infty$ and $-1$ as $k\to\infty$. Thus the root exists uniquely.

The radius of $\beta_N$ is positive. If two positions agree, their radii
give equal $\sin Nt$ values, after which their arguments give equal
parameters modulo $2\pi$. Thus $\beta_N$ is embedded. With
$c=\cos Nt$ and $s=\sin Nt$,
\begin{equation}\label{eq:corrugated-derivative}
 \beta_N'=e^{i(t+2\epsilon s)}
       \bigl(c+i(1+\epsilon s)(1+2c)\bigr).
\end{equation}
The real and imaginary components cannot vanish together, proving
regularity. Its angle lift increases by $2\pi$ over the period, proving
positive orientation and enclosure of zero.

The derivative $\delta_N'=m_N\beta_N'$ is covariant because $m_N$ is
cell-periodic. The chosen initial value gives
$\delta_N(T_N)=\xi_N\delta_N(0)$, hence covariance for the whole curve
and closure after $N$ cells. Finally
\[
 \det(\beta_N,\delta_N')
 =m_N(1+\epsilon\sin Nt)^2(1+2\cos Nt),
\]
so the scalar root cancels the mixed action on each cell and on the full
period.

\subsection{Position and tangent bounds}
Put
\[
 L(c)=\sqrt{c^2+(1+2c)^2},\qquad
 U_0=e^{it}\frac{c+i(1+2c)}{L(c)},\qquad
 U_N=\frac{\beta_N'}{|\beta_N'|}.
\]
For $N\ge10$ we claim
\begin{equation}\label{eq:corrugated-estimates}
 |\beta_N-e^{it}|\le3/N,\qquad
 |U_N-U_0|<50/N,\qquad
 |\delta_N-(i/2)e^{it}|<60/N.
\end{equation}
The first uses $|e^{iu}-1|\le|u|$. Since
$1/\sqrt5\le L(c)\le\sqrt{10}$, the unnormalized tangent in
\eqref{eq:corrugated-derivative} differs from
$e^{it}(c+i(1+2c))$ by less than $10/N$. The bound
\[
 \left|\frac z{|z|}-\frac w{|w|}\right|\le\frac{2|z-w|}{|w|}
\]
therefore gives the tangent estimate.

To control the partner, use the normalized weighted average
\[
 \langle f\rangle_N=\frac1{2\pi}\int_0^{2\pi}
                Z_N^{-1}e^{-k_N\cos x}f(x)\,dx.
\]
Odd functions of $\sin x$ have zero average. With $y=1+2\cos x$,
the root identity gives
\[
 \langle y\rangle_N=-\epsilon^2\langle\sin^2x\,y\rangle_N,
 \quad |\langle y\rangle_N|\le3\epsilon^2,
 \quad\left|\langle\cos x+iy\rangle_N+\tfrac12\right|
          \le\tfrac92\epsilon^2.
\]
The one-cell average
\[
 J_N=\left\langle e^{i\epsilon(x+2\sin x)}
       \bigl(\cos x+i(1+\epsilon\sin x)y\bigr)\right\rangle_N
\]
satisfies $|J_N+1/2|<25\epsilon$: use
$0\le x+2\sin x\le2\pi$ and the preceding average bounds.
For $\alpha=2\pi\epsilon$, write $Q(\alpha)=\alpha/(e^{i\alpha}-1)$.
For $\epsilon\le1/10$, $|Q|<11/10$ and $|Q+i|<4\epsilon$.
Since $\delta_N(0)=Q(\alpha)J_N$, it follows that
$|\delta_N(0)-i/2|<30\epsilon$.
Over one cell,
\[
 |\delta_N(t)-\delta_N(0)|
 \le2\pi\epsilon(\sqrt{10}+3\epsilon)<22\epsilon,
 \qquad |(i/2)e^{it}-i/2|<4\epsilon.
\]
Their sum is less than $60\epsilon$. Covariance extends this estimate
to the full period, completing \eqref{eq:corrugated-estimates}.

\subsection{The two visibility margins}
At $\delta_0=(i/2)e^{it}$,
\[
 V_{1/4}(\delta_0)=e^{it}(-\sqrt3/2+i/2).
\]
Thus
\begin{equation}\label{eq:outer-limit-visibility}
 U_0\cdot V_{1/4}(\delta_0)
 =\frac{1+(2-\sqrt3)c}{2L(c)}
 \ge\frac{\sqrt3-1}{2\sqrt{10}}>\frac1{10}.
\end{equation}
For the reflected reversed pair, the identity $A_*J=-JA_*$ reduces
visibility to positivity of
\[
 U_N\cdot\frac{-R_-\beta_N+
             \sqrt{|\beta_N|^2-R_-^2}\,J\beta_N}{|\beta_N|^2}.
\]
At $\beta_0=e^{it}$ and $R_-=4/5$ the limiting expression is
\begin{equation}\label{eq:inner-limit-visibility}
 \frac{-4c+3(1+2c)}{5L(c)}
 =\frac{3+2c}{5L(c)}\ge\frac1{5\sqrt{10}}>\frac1{16}.
\end{equation}
The differential norm of either direction map
$z\mapsto(\pm Rz+\sqrt{|z|^2-R^2}Jz)/|z|^2$ is
$1/\sqrt{|z|^2-R^2}$. For $N\ge10^4$, the comparison segments in
\eqref{eq:corrugated-estimates} have radii at least $0.494$ for the
outer test and $0.9997$ for the inner test. A Lipschitz bound of $3$
therefore suffices. The respective errors are at most
$50/N+3(60/N)=230/N$ and $50/N+3(3/N)=59/N$.
Both are smaller than their displayed margins, and both radius
inequalities hold. Since the velocities in each pair are positive
multiples, both required tangent inequalities follow. The positive-partner
lemma proves embedding and enclosure for $\delta_N$.

\subsection{The spherical curvature is nonconstant}
The gnomonic lift is embedded and stays in the northern hemisphere.
With the sign convention \eqref{eq:spherical-frame}, its geodesic
curvature has the sign of $\det(\beta_N',\beta_N'')$. Writing
$r=1+\epsilon\sin Nt$, $c=\cos Nt$, and $s=\sin Nt$, differentiation gives
\[
 \det(\beta_N',\beta_N'')
 =Nr s+2c^2(1+2c)+r^2(1+2c)^3.
\]
The values at $Nt=0$ and $Nt=\pi$ are $33$ and $-3$.
The curvature therefore takes both signs. This finishes every assertion
of Lemma~\ref{lem:seed}.

\begin{thebibliography}{99}
\bibitem{GhomiProblems}
M.~Ghomi, ``Open problems in geometry of curves and surfaces,'' 2019,
Problems 1.1--1.4,
\href{https://ghomi.math.gatech.edu/Papers/op.pdf}{author's survey}.

\bibitem{Kovaleva}
G.~A.~Kovaleva, ``Example of a surface that is homotopic to a tube and has a
closed asymptotic line,'' \emph{Mathematical Notes} \textbf{3} (1968),
257--263, \href{https://www.mathnet.ru/eng/mzm6695}{MathNet}.

\bibitem{GarciaSotomayor}
R.~Garcia and J.~Sotomayor, ``Structural stability of parabolic points and
periodic asymptotic lines,'' \emph{Matem\'atica Contempor\^anea}
\textbf{12} (1997), 83--102,
\href{https://repositorio.bc.ufg.br/items/02c6dfb8-0487-4267-a20d-4d97c54decf0}{repository}.

\bibitem{DaCruzGarcia}
D.~H.~da Cruz and R.~A.~Garcia, ``Finite type $\xi$-asymptotic lines of plane
fields in $\mathbb R^3$,'' \emph{Journal of Singularities} \textbf{22} (2020),
17--27, \href{https://doi.org/10.5427/jsing.2020.22b}{doi:10.5427/jsing.2020.22b}.

\bibitem{SaEarpToubiana}
R.~Sa Earp and E.~Toubiana, ``Discrete and nondiscrete isometric deformations
of surfaces in $\mathbb R^3$,'' \emph{Siberian Mathematical Journal}
\textbf{43} (2002), 714--718,
\href{https://doi.org/10.1023/A:1016384521524}{doi:10.1023/A:1016384521524};
\href{https://m.mathnet.ru/eng/smj1338}{MathNet}.

\bibitem{Sabitov2018}
I.~Kh.~Sabitov, ``On a problem in the bendings theory of negatively curved
surfaces,'' \emph{Siberian Electronic Mathematical Reports} \textbf{15} (2018),
890--893,
\href{https://doi.org/10.17377/semi.2018.15.076}{doi:10.17377/semi.2018.15.076}.

\bibitem{Sabitov2021}
I.~Kh.~Sabitov, ``On metrics admitting a discontinuum set of non-congruent
immersions in $\mathbb R^3$,'' \emph{Siberian Electronic Mathematical Reports}
\textbf{18} (2021), 1023--1026,
\href{https://www.mathnet.ru/eng/semr1418}{MathNet}.

\bibitem{Ding2026}
W.~Ding, ``A solution to a problem of Ghomi and Raffaelli on closed
asymptotic curves,'' 2026, preprint (version~1, October~6),
\href{https://arxiv.org/abs/2610.08050}{arXiv:2610.08050}.

\bibitem{NiZhangZhang}
L.~Ni, W.~Zhang, and Y.~Zhang, ``On questions of Pogorelov and Toponogov,''
2026, preprint (version~1, June~28),
\href{https://arxiv.org/abs/2606.29231}{arXiv:2606.29231}.

\bibitem{BK}
T. F. Banchoff and W. K\"uhnel,
\emph{Tight submanifolds, smooth and polyhedral}, in T. E. Cecil and S.-S. Chern (eds.), \emph{Tight and Taut Submanifolds}, MSRI Publications \textbf{32}, Cambridge University Press, 1997, 51--118.
\href{https://doi.org/10.1017/9781009701815.003}{doi:10.1017/9781009701815.003}.

\bibitem{Nirenberg}
L. Nirenberg,
\emph{Rigidity of a class of closed surfaces}, in \emph{Nonlinear Problems (Proc. Sympos., Madison, Wis., 1962)}, University of Wisconsin Press, Madison, 1963, 177--193.

\bibitem{YauReview}
S.-T. Yau, ``Review of geometry and analysis,'' \emph{Asian Journal of
Mathematics} \textbf{4} (2000), no.~1, 235--278,
\href{https://archive.intlpress.com/site/pub/files/_fulltext/journals/ajm/2000/0004/0001/AJM-2000-0004-0001-a016.pdf}{publisher PDF}.

\bibitem{GR}
M. Ghomi and M. Raffaelli,
\emph{Topology of closed asymptotic curves on negatively curved surfaces}, Journal of Geometric Analysis \textbf{35} (2025), article 381.
\href{https://doi.org/10.1007/s12220-025-02200-3}{doi:10.1007/s12220-025-02200-3};
\href{https://arxiv.org/abs/2412.19266}{arXiv:2412.19266}.

\bibitem{HK}
Q. Han and M. Khuri,
\emph{Rigidity in the class of orientable compact surfaces of minimal total absolute curvature}, Differential Geometry and its Applications \textbf{29} (2011), 463--472.
\href{https://doi.org/10.1016/j.difgeo.2011.04.035}{doi:10.1016/j.difgeo.2011.04.035};
\href{https://arxiv.org/abs/1104.0474}{arXiv:1104.0474}.

\bibitem{IzmProj}
I. Izmestiev,
\emph{Projective background of the infinitesimal rigidity of frameworks}, arXiv:0804.2694 (2008), revised 2009.
\href{https://doi.org/10.48550/arXiv.0804.2694}{doi:10.48550/arXiv.0804.2694}.

\bibitem{Rolfsen}
D.~Rolfsen, \emph{Knots and Links}, AMS Chelsea Publishing \textbf{346}, American Mathematical Society, Providence, RI, 2003; reprint of the 1976 original.
\end{thebibliography}

\end{document}
